\documentclass[11pt,a4paper,spanish]{book}
\usepackage{estilo_unir-1}
\hypersetup{
	colorlinks   = true,
	urlcolor     = blue,
	linkcolor    = black,
	citecolor   = red,
	pdftitle     = {Marco Híbrido GNN-HVA para la Caracterización de Fases Cuánticas Orientado a Hardware NISQ},
	pdfauthor    = {Franco Raineri},
	pdfsubject   = {Trabajo de Fin de Máster - UNIR},
	pdfkeywords  = {VQE, GNN, transiciones de fase cuánticas, Ising, circuitos variacionales, HVA}
}
\usepackage[nottoc]{tocbibind}
\usepackage{booktabs}
\usepackage{multirow}
\usepackage{siunitx}

%---------------------------
%título de la asignatura, trabajo y autor
%---------------------------
\subject{Trabajo de Fin de Máster}
\title{Marco Híbrido GNN-HVA para la Caracterización de Fases Cuánticas Orientado a Hardware NISQ}
\author{Franco Raineri}
\profesor{Andrés Navas Cáliz}
\date{Agosto de 2026}

%---------------------------
\begin{document}
\renewcommand{\listfigurename}{Índice de figuras}
\renewcommand{\listtablename}{Índice de tablas}
\renewcommand{\contentsname}{Índice de contenidos}
\renewcommand{\figurename}{Figura}
\renewcommand{\tablename}{Tabla}

\maketitle
\frontmatter
\tableofcontents
\listoffigures
\listoftables

\chapter{Resumen}

El presente Trabajo Fin de Máster desarrolla y valida un marco híbrido clásico-cuántico para la caracterización de fases cuánticas, integrando técnicas de predicción de parámetros mediante redes neuronales de grafos con algoritmos variacionales cuánticos. El marco combina una Red Neuronal de Grafos de Paso de Mensajes (MPNN) basada en GINConv con un Ansatz Variacional Hamiltoniano (HVA), operando sobre el Modelo de Ising en Campo Transversal (TFIM) como sistema de validación. Siguiendo el paradigma de predicción de parámetros propuesto en la literatura reciente, la MPNN predice los parámetros óptimos del circuito variacional directamente desde la representación en grafo del Hamiltoniano, eliminando el bucle de optimización iterativa del VQE. La contribución principal es la integración de estas técnicas en un pipeline unificado y su validación sistemática sobre 400+ ejecuciones en simulación ideal, abarcando 5 topologías de red (cadena 1D, escalera, triangular, cuadrada, heavy-hex), tamaños de sistema desde $N=4$ hasta $N=20$ qubits, y profundidades de circuito $p = 1$--$4$, demostrando que el pipeline mantiene $\Delta E/\text{gap} < 5\%$ de manera consistente dentro del régimen operativo válido, con una aceleración de $29$--$500\times$ respecto a VQE con inicialización aleatoria. Se valida además la extensibilidad del marco a variantes del Hamiltoniano de Ising (campo longitudinal, modelo frustrado) y se documentan formalmente los límites de expresividad del ansatz (Heisenberg, Kitaev). Los resultados establecen una metodología reproducible para la clasificación de fases cuánticas mediante observables locales, sentando las bases para su despliegue futuro en hardware cuántico real.

\vspace{0.5cm}
{\bf Palabras clave:} Computación cuántica, VQE, Redes Neuronales de Grafos, transiciones de fase cuánticas, Modelo de Ising, circuitos variacionales, Ansatz Variacional Hamiltoniano.

\chapter{Abstract}

This Master's Thesis develops and validates a hybrid classical-quantum framework for quantum phase characterization, integrating GNN-based parameter prediction with variational quantum circuits into a unified pipeline. The framework combines a Message-Passing Neural Network (MPNN) based on GINConv with a Hamiltonian Variational Ansatz (HVA), operating on the Transverse Field Ising Model (TFIM) as the validation system. Following the parameter prediction paradigm from recent literature, the MPNN predicts optimal variational circuit parameters directly from the Hamiltonian's graph representation, eliminating the iterative optimization loop of VQE. The main contribution is the integration of these techniques and their systematic validation over 400+ ideal-simulation pipeline executions, spanning 5 lattice topologies (chain 1D, ladder, triangular, square, heavy-hex), system sizes from $N=4$ to $N=20$ qubits, and circuit depths $p = 1$--$4$, demonstrating that the pipeline maintains $\Delta E/\text{gap} < 5\%$ consistently within the valid operating regime, with a $29$--$500\times$ speedup over randomly-initialized VQE. Framework extensibility to Ising Hamiltonian variants (longitudinal field, frustration) is validated, and ansatz expressibility limits are formally documented (Heisenberg, Kitaev). The results establish a reproducible methodology for quantum phase classification via local observables, laying the groundwork for future deployment on quantum hardware.

\vspace{0.5cm}
{\bf Keywords:} Quantum computing, VQE, Graph Neural Networks, quantum phase transitions, Ising model, variational circuits, Hamiltonian Variational Ansatz.

\mainmatter

%---------------------------
% CAPÍTULO 1: INTRODUCCIÓN
%---------------------------

\chapter{Introducción}

La caracterización de transiciones de fase cuánticas constituye uno de los problemas centrales de la física de la materia condensada. Cuando un sistema de muchos cuerpos experimenta cambios cualitativos en su estado fundamental al variar un parámetro externo ---sin intervención de fluctuaciones térmicas--- emergen fenómenos de entrelazamiento de largo alcance que hacen intratable la simulación clásica \citep{dutta2015, anderson1972}. El espacio de Hilbert crece exponencialmente con el número de partículas ($2^N$ para espines-1/2), y los métodos estadísticos como Monte Carlo Cuántico fracasan ante el Problema del Signo para sistemas frustrados \citep{troyer2005}.

La simulación cuántica, propuesta por \citet{feynman1982}, ofrece una solución natural. Los algoritmos variacionales como VQE \citep{peruzzo2014} representan la aproximación más prometedora para hardware de escala intermedia, pero enfrentan tres barreras: (1) barren plateaus en circuitos profundos \citep{mcclean2018, cerezo2021}, (2) truncamiento por ruido que limita la profundidad efectiva a $O(\log n)$ capas \citep{mele2026}, y (3) costo de evaluación prohibitivo (cientos a miles de evaluaciones por punto \citep{tilly2022}).

La convergencia de tres resultados teóricos recientes abre una ventana de operación viable: las funciones de costo locales están libres de barren plateaus \citep{cerezo2021}, el HVA maximiza expresividad por parámetro con paisaje suave dentro de una fase \citep{wiersema2020, mele2022, puig2025}, y las GNN pueden aprender propiedades del estado fundamental generalizando a nuevas instancias \citep{huang2022, miao2024, zhang2025}.

\section{Planteamiento del problema y objetivo}

Este Trabajo propone la integración de predicción de parámetros mediante GNN \citep{miao2024, zhang2025} y optimización warm-start \citep{puig2025} en una \textbf{Arquitectura Predictiva Híbrida (GNN-HVA)}, implementada como un pipeline de tres fases:

\begin{enumerate}
    \item \textbf{Fase 1 --- Ground Truth:} Diagonalización exacta / DMRG \citep{schollwock2011}.
    \item \textbf{Fase 2 --- VQE Warm-Start:} Optimización del HVA mediante barrido descendente, donde cada punto hereda los parámetros óptimos del anterior \citep{puig2025, mele2022}.
    \item \textbf{Fase 3 --- Predictor MPNN:} GINConv \citep{xu2019} predice parámetros óptimos en una pasada clásica; evaluación mediante simulación ideal.
\end{enumerate}

El pipeline se valida íntegramente mediante simulación ideal (\texttt{StatevectorEstimator}), aislando la contribución del framework de los efectos de ruido. El despliegue en hardware cuántico real con mitigación de errores \citep{kim2023} constituye trabajo futuro.

\textbf{Objetivo general:} Demostrar que la integración de predicción GNN con HVA en un pipeline unificado permite reducir el costo de la clasificación de fases entre $29\times$ y $500\times$ respecto a VQE convencional, manteniendo $\Delta E/\text{gap} < 5\%$ dentro del régimen operativo válido, y documentar formalmente los límites de dicho régimen.

Como sistema de validación se emplea el TFIM \citep{dutta2015}, con extensiones a campo longitudinal y modelo frustrado. Se documenta el límite de expresividad para Heisenberg. La validación abarca 5 topologías, profundidades $p = 1$--$4$, y tamaños $N = 4$--$20$ (400+ ejecuciones).

\textbf{Posicionamiento y relevancia para la materia condensada.} A las escalas de validación empleadas ($N = 4$--$20$, 1D), los resultados son reproducibles mediante simuladores clásicos exactos ---la frontera de ventaja cuántica se sitúa en $N \approx 20$ para sistemas 2D \citep{martin2026}. El valor del presente trabajo es metodológico: se valida el pipeline en un régimen donde la respuesta exacta es conocida y verificable, estableciendo la infraestructura para su aplicación futura en sistemas donde los métodos clásicos son intratables (redes frustradas 2D, Kagomé a $N > 100$ \citep{ahsan2025}). En el contexto de la investigación experimental en materia condensada, el pipeline ofrece tres utilidades concretas: (i) screening rápido de diagramas de fases en procesadores cuánticos, donde un entrenamiento inicial con 17--27 puntos VQE permite predecir el resto del espacio de parámetros sin optimización adicional; (ii) clasificación de fases mediante observables locales ($\langle X_i \rangle$, $\langle Z_i Z_{i+1} \rangle$) sin reconstrucción de estado, compatible con las limitaciones de medición del hardware ruidoso \citep{sharma2026}; y (iii) establecimiento de una metodología de benchmark reproducible para evaluar ansätze y predictores en el contexto de la simulación cuántica variacional.

\section{Estructura del Trabajo}

El resto de la presente memoria se estructura en cinco capítulos principales:

\begin{itemize}
    \item \textbf{Capítulo 2: Marco Teórico y Estado del Arte.} Presenta los fundamentos físicos y matemáticos: el TFIM, VQE, HVA, truncamiento por ruido \citep{mele2026}, barren plateaus y redes neuronales de grafos. Se analiza críticamente el estado del arte en predicción de parámetros variacionales, identificando la brecha que este trabajo pretende llenar.

    \item \textbf{Capítulo 3: Desarrollo del Trabajo.} Describe la metodología completa del pipeline de tres fases, incluyendo la arquitectura del software, los algoritmos de entrenamiento, las estrategias de validación y los criterios de éxito. Detalla la implementación de las variantes del Hamiltoniano de Ising empleadas para las pruebas sistemáticas.

    \item \textbf{Capítulo 4: Descripción Detallada e Identificación de Requisitos.} Detalla el problema a resolver en su contexto computacional y físico, establece los requisitos funcionales del pipeline, y describe el trabajo previo realizado para delimitar el alcance de la solución.

    \item \textbf{Capítulo 5: Evaluación.} Presenta todos los resultados numéricos obtenidos mediante simulación ideal: validación por topología y escala del sistema, validación de escalamiento con la profundidad del circuito, extensibilidad a variantes del Hamiltoniano de Ising, comparativa cuantitativa con procedimientos de la literatura, y análisis de límites físicos del ansatz. Incluye la sección de código desarrollado con referencia al repositorio.

    \item \textbf{Capítulo 6: Conclusiones y Trabajo Futuro.} Sintetiza las contribuciones, evalúa el cumplimiento de los objetivos, resume la ventaja de la solución propuesta, y propone líneas de investigación futuras ---incluyendo el despliegue en hardware cuántico real con mitigación de errores.
\end{itemize}


%---------------------------
% CAPÍTULO 2: MARCO TEÓRICO Y ESTADO DEL ARTE
%---------------------------

\chapter{Marco Teórico y Estado del Arte}

Este capítulo presenta los fundamentos teóricos en las tres disciplinas que convergen en este trabajo: la física de transiciones de fase cuánticas, los algoritmos cuánticos variacionales y las redes neuronales de grafos. El objetivo es establecer el marco conceptual necesario para comprender la arquitectura híbrida propuesta e identificar las limitaciones del estado del arte actual.

\section{El problema cuántico de muchos cuerpos}

El Problema de los Muchos Cuerpos Cuántico (Quantum Many-Body Problem) constituye uno de los desafíos computacionales fundamentales de la física moderna. Dado un sistema de $N$ partículas cuánticas interactuantes, su estado se describe mediante un vector en un espacio de Hilbert cuya dimensión crece exponencialmente como $2^N$ para espines-1/2. Para $N = 40$ partículas, el espacio de estados alcanza $\sim 10^{12}$ dimensiones, superando la capacidad de cualquier supercomputador clásico existente \citep{feynman1982}.

\citet{troyer2005} demostraron que la simulación general de sistemas cuánticos fermiónicos mediante Monte Carlo es NP-hard debido al Problema del Signo: las amplitudes de probabilidad pueden ser negativas, causando cancelaciones que impiden la convergencia estadística. Este resultado establece que, salvo $P = NP$, no existe algoritmo clásico eficiente para simular la dinámica general de estos sistemas.

\subsection{Transiciones de Fase Cuánticas}

A diferencia de las transiciones de fase clásicas (gobernadas por fluctuaciones térmicas), las transiciones de fase cuánticas (QPT) ocurren a temperatura cero y están impulsadas por fluctuaciones puramente cuánticas al variar un parámetro del Hamiltoniano \citep{dutta2015, sachdev2011}. En el punto crítico cuántico, las correlaciones divergen y el sistema exhibe entrelazamiento de largo alcance, haciendo que los métodos de campo medio y las aproximaciones perturbativas fallen simultáneamente.

La clasificación de fases cuánticas se realiza mediante parámetros de orden: cantidades observables que toman valores distintos en cada fase. Para las transiciones de ruptura de simetría ---como la del Modelo de Ising--- estos parámetros de orden pueden expresarse mediante observables locales, lo que resulta crucial para su medición en hardware cuántico ruidoso.

\subsection{El Modelo de Ising en Campo Transversal (TFIM)}

El TFIM es el modelo paradigmático para el estudio de transiciones de fase cuánticas \citep{dutta2015}. Su Hamiltoniano para una cadena unidimensional de $N$ espines con condiciones de contorno abiertas se escribe como:

\begin{equation}
H_{\text{TFIM}} = -J \sum_{i=1}^{N-1} Z_i Z_{i+1} - h \sum_{i=1}^{N} X_i
\label{eq:tfim}
\end{equation}

\noindent donde $J > 0$ es la constante de acoplamiento entre espines vecinos, $h$ es la intensidad del campo magnético transversal, y $Z_i$, $X_i$ son las matrices de Pauli actuando sobre el sitio $i$.

El modelo exhibe dos fases bien diferenciadas:

\begin{itemize}
    \item \textbf{Fase ferromagnética} ($h/J \ll 1$): Los espines se alinean a lo largo del eje $z$, rompiendo espontáneamente la simetría $\mathbb{Z}_2$. El parámetro de orden es $\langle Z_i \rangle \neq 0$.
    \item \textbf{Fase paramagnética} ($h/J \gg 1$): El campo transversal domina y los espines se alinean con el campo. La magnetización transversal $\langle X_i \rangle \to 1$.
\end{itemize}

En el límite termodinámico ($N \to \infty$), la transición ocurre en $h_c/J = 1.0$. Para sistemas finitos, la transición se suaviza en una región de crossover cuyo ancho escala como $\sim N^{-1/\nu}$ con $\nu = 1$ para la clase de universalidad del TFIM \citep{dutta2015, sachdev2011}.

En este trabajo se consideran además extensiones del Hamiltoniano de Ising que permiten validar la generalidad del pipeline:

\begin{equation}
H_{\text{TFIM+long}} = H_{\text{TFIM}} - g \sum_{i=1}^{N} Z_i
\label{eq:tfim_long}
\end{equation}

\noindent que añade un campo longitudinal $g$ sin incrementar el número de compuertas de dos qubits del circuito, y

\begin{equation}
H_{\text{J1-J2}} = -J_1 \sum_{\langle i,j \rangle} Z_i Z_j - J_2 \sum_{\langle\langle i,j \rangle\rangle} Z_i Z_j - h \sum_{i=1}^{N} X_i
\label{eq:tfim_frustrated}
\end{equation}

\noindent que introduce frustración geométrica mediante acoplamientos a segundos vecinos $J_2$.

La elección del TFIM como sistema de validación se justifica por tres razones:

\begin{enumerate}
    \item \textbf{Solución exacta conocida:} Permite verificación rigurosa del pipeline contra resultados analíticos.
    \item \textbf{Mapeo isomórfico a qubits:} Cada espín-1/2 corresponde directamente a un qubit, sin overhead de codificación.
    \item \textbf{Observables locales suficientes:} La transición se caracteriza mediante $\langle X_i \rangle$ y $\langle Z_i Z_{i+1} \rangle$ \citep{dutta2015}, medibles con presupuestos polinomiales de shots en hardware cuántico \citep{cerezo2021}.
\end{enumerate}

\section{Algoritmos Cuánticos Variacionales}

\subsection{El Eigensolver Cuántico Variacional (VQE)}

El VQE, propuesto por \citet{peruzzo2014}, es un algoritmo híbrido clásico-cuántico basado en el principio variacional:

\begin{equation}
E_0 \leq \langle \psi(\boldsymbol{\theta}) | H | \psi(\boldsymbol{\theta}) \rangle = E(\boldsymbol{\theta})
\label{eq:variational}
\end{equation}

\noindent donde $|\psi(\boldsymbol{\theta})\rangle = U(\boldsymbol{\theta})|0\rangle$ es un estado preparado por un circuito cuántico parametrizado. La revisión de \citet{tilly2022} establece que la eficacia de VQE depende críticamente de la elección del ansatz, la función de costo y la estrategia de inicialización de parámetros.

\subsection{El Ansatz Variacional Hamiltoniano (HVA)}

El HVA, introducido por \citet{wiersema2020}, construye el circuito parametrizado a partir de la estructura del Hamiltoniano. Para el TFIM, el HVA de $p$ capas toma la forma:

\begin{equation}
|\psi(\boldsymbol{\theta})\rangle = \prod_{l=1}^{p} \left[ e^{-i\theta_{l,x} \sum_i X_i} \cdot e^{-i\theta_{l,zz} \sum_{\langle i,j \rangle} Z_i Z_j} \right] |+\rangle^{\otimes N}
\label{eq:hva}
\end{equation}

\noindent donde $|+\rangle^{\otimes N}$ es el estado inicial (producto de autoestados de $X$), que corresponde al estado fundamental exacto en el límite $h \to \infty$.

Las ventajas del HVA frente a ansätze genéricos tipo Hardware-Efficient (HEA) han sido confirmadas por \citet{tripathi2026} en un estudio comparativo sobre TFIM en 1D, 2D y 3D hasta 27 espines:

\begin{itemize}
    \item \textbf{Respeto de simetrías:} El HVA preserva las simetrías del Hamiltoniano por construcción.
    \item \textbf{Expresividad por parámetro:} Con solo $2p$ parámetros alcanza fidelidades superiores al HEA ($O(Np)$ parámetros).
    \item \textbf{Paisaje suave:} La transferibilidad de parámetros genera un paisaje sin mínimos locales espurios \citep{mele2022}.
    \item \textbf{Ausencia de barren plateaus:} Para profundidades acotadas, exhibe gradientes polinomialmente decrecientes \citep{wiersema2020}.
\end{itemize}

\subsection{Barren Plateaus}

El fenómeno de barren plateaus (BP) fue identificado por \citet{mcclean2018}:

\begin{equation}
\text{Var}\left[\frac{\partial C}{\partial \theta_k}\right] \leq F(n) \in O(2^{-n})
\label{eq:bp}
\end{equation}

\citet{cerezo2021} demostraron que para funciones de costo locales y circuitos de profundidad $O(\log n)$:

\begin{equation}
\text{Var}\left[\frac{\partial C_{\text{local}}}{\partial \theta_k}\right] \geq \Omega\left(\frac{1}{\text{poly}(n)}\right)
\label{eq:local_bp}
\end{equation}

\noindent Esta cota inferior polinomial garantiza gradientes detectables, fundamentando la elección de observables locales como función de costo.

\subsection{Truncamiento de Profundidad por Ruido No Unital}

\citet{mele2026} demostraron en Nature Physics que los canales de ruido no unital ---incluyendo relajación térmica y amortiguamiento de amplitud, predominantes en hardware superconductor--- poseen un punto fijo único $\rho_{\text{fijo}}$. Para un circuito de profundidad $d$ con tasa de ruido $\gamma$ por capa, el estado de salida $\rho_d$ satisface:

\begin{equation}
\|\rho_d - \rho_{\text{fijo}}\|_1 \leq (1 - \gamma)^d
\label{eq:truncation}
\end{equation}

\noindent La profundidad efectiva ---las capas que contribuyen de forma detectable al valor esperado de un observable--- está acotada por:

\begin{equation}
d_{\text{eff}} = O\!\left(\frac{\log(n/\varepsilon)}{\gamma}\right)
\label{eq:depth_bound}
\end{equation}

\noindent donde $n$ es el número de qubits y $\varepsilon$ la precisión deseada en la estimación del observable. Las capas más allá de $d_{\text{eff}}$ contribuyen una señal exponencialmente suprimida: la información que codifican se disipa irrecuperablemente en el punto fijo del ruido. Crucialmente, los autores demuestran el resultado complementario positivo: en el régimen $d \leq d_{\text{eff}}$, las funciones de costo locales \emph{no} sufren barren plateaus, garantizando gradientes detectables para la optimización.

\textbf{Implicación para la era NISQ.} La Ecuación~\ref{eq:depth_bound} establece que la profundidad del circuito variacional debe mantenerse acotada en función de la tasa de ruido del hardware. No impone un valor fijo de $p$, sino una relación inversa entre profundidad útil y ruido: a mayor $\gamma$, menor es el número de capas que aportan información. Para hardware contemporáneo (IBM Heron, tasas de error CZ $\sim 0.3$--$0.5\%$), la cota permite decenas de capas en principio, pero cada capa adicional incrementa el presupuesto total de puertas de dos qubits y, con ello, el costo de mitigación de errores ---que escala exponencialmente con la profundidad del circuito y el número de qubits \citep{tsubouchi2023}. En consecuencia, la profundidad operativa óptima resulta de equilibrar expresividad del ansatz contra costo de mitigación.

Esto establece la sinergia teórica que fundamenta la arquitectura del pipeline:
\begin{enumerate}
    \item \textbf{Profundidad acotada} $\to$ el circuito opera dentro de $d_{\text{eff}}$, preservando la señal frente al ruido.
    \item \textbf{Funciones de costo locales} $\to$ gradientes detectables, sin barren plateaus \citep{cerezo2021}.
    \item \textbf{Estructura HVA} $\to$ máxima expresividad por parámetro en el régimen de profundidad acotada \citep{wiersema2020}, con paisaje energético suave que habilita transferibilidad \citep{mele2022}.
\end{enumerate}

\noindent Esta triple convergencia ---profundidad acotada, costos locales y estructura HVA--- define el régimen donde los algoritmos variacionales son simultáneamente entrenables, resilientes al ruido y físicamente expresivos. En principio, la validación empírica de \citet{tripathi2026} confirma que el HVA con profundidad moderada es suficiente para clasificar correctamente las fases del TFIM hasta $N = 27$ qubits.

\section{Estrategias de Inicialización y Warm-Start}

\citet{mele2022} demostraron que los parámetros óptimos del HVA varían suavemente con los parámetros del Hamiltoniano, fundamentando la estrategia de warm-start por barrido descendente. \citet{puig2025} formalizaron este resultado en PRX Quantum, demostrando que el warm-start proporciona varianzas de pérdida mayores (señales de gradiente más fuertes) comparadas con la inicialización aleatoria. \citet{schiffer2026} derivaron recientemente las condiciones bajo las cuales la optimización basada en gradientes con warm-start adiabático mantiene exitosamente el seguimiento del estado fundamental a lo largo de una secuencia de Hamiltonianos, proporcionando garantías formales de convergencia para nuestra estrategia de barrido descendente. \citet{skogh2023} validaron experimentalmente la transferencia de parámetros entre instancias cercanas del Hamiltoniano.

La extensión natural del warm-start secuencial es la predicción directa mediante modelos de aprendizaje automático. A partir de 2024, este paradigma ha sido validado por múltiples grupos de manera independiente, consolidándose como una línea de investigación activa con al menos siete trabajos publicados:

\textbf{NN-VQE} \citep{miao2024}: Un perceptrón multicapa entrenado con 20 puntos de datos predice parámetros VQE óptimos para Hamiltonianos de espines parametrizados, empleando dropout y aprendizaje activo. Primer trabajo en demostrar que un modelo clásico puede reemplazar el bucle de optimización VQE.

\textbf{Qracle} \citep{zhang2025}: Un inicializador GNN que codifica conjuntamente el Hamiltoniano y el ansatz como grafo unificado, reduciendo hasta un 64\% las iteraciones de optimización VQE. Introduce la representación en grafo del circuito variacional como entrada al predictor.

\textbf{Flow-VQE} \citep{zou2026}: Flujos normalizantes generativos para warm-start de VQE, alcanzando aceleraciones de $50\times$. Emplea un enfoque probabilístico que captura la multimodalidad del paisaje energético.

\textbf{PVLS} \citep{yang2024pvls}: Predicción de parámetros para solucionadores cuánticos variacionales lineales (VQLS), extendiendo el paradigma de inicialización ML más allá de VQE hacia otros algoritmos variacionales.

\textbf{VQEzy} \citep{zhang2025vqezy}: Dataset abierto de parámetros VQE optimizados para múltiples Hamiltonianos y ansätze, proporcionando un benchmark estandarizado para evaluar predictores de parámetros. Evidencia la madurez del campo al crear infraestructura compartida.

\textbf{PALQO} \citep{huang2025palqo}: Modelo informado por la física para optimización cuántica a gran escala, publicado en NeurIPS 2025. Emplea simetrías del Hamiltoniano como sesgo inductivo del predictor, alcanzando generalización a sistemas de hasta 100 qubits.

\textbf{Fast ML-VQE} \citep{karim2025}: Optimizador VQE basado en ML que incorpora modelos de ruido del hardware en el entrenamiento, produciendo predicciones directamente adaptadas a dispositivos ruidosos sin necesidad de mitigación posterior.

Adicionalmente, \citet{lee2026} demostraron que un autoencoder de grafos puede generar parámetros VQE que generalizan entre instancias de Hamiltonianos sin optimización específica.

La proliferación de estos trabajos confirma que la predicción de parámetros variacionales mediante ML es un paradigma consolidado. Sin embargo, cada trabajo aborda un aspecto específico ---arquitectura del predictor, tipo de algoritmo variacional, incorporación de ruido, o generación de datasets--- sin que ninguno integre la validación sistemática sobre múltiples topologías de red, profundidades de circuito y variantes del Hamiltoniano en un pipeline unificado reproducible.

\section{Redes Neuronales de Grafos}

\subsection{Fundamentos del Paso de Mensajes}

Las Redes Neuronales de Grafos (GNN) operan sobre datos estructurados como grafos $\mathcal{G} = (V, E)$. El paradigma de Paso de Mensajes, formalizado por \citet{gilmer2017}, define la actualización iterativa de representaciones nodales mediante agregación de información de vecinos. Para sistemas de espines, la representación como grafo es natural e isomórfica: qubits son nodos, interacciones son aristas, y campos locales son atributos de nodo.

\subsection{GINConv: Máxima Expresividad para Grafos}

\citet{xu2019} demostraron que la Graph Isomorphism Network (GIN) es tan expresiva como el test de Weisfeiler-Lehman para distinguir grafos no isomorfos ---la cota superior teórica para MPNNs. Para redes uniformes como la cadena 1D del TFIM, donde todas las aristas son equivalentes, mecanismos de atención (GATConv) añaden parámetros sin ganancia informativa. La superioridad de GNN sobre CNN (+36\%) para predicción de propiedades de circuitos cuánticos fue confirmada por \citet{meng2025}. \citet{slavin2025} validaron independientemente que la geometría de la red es suficiente para predecir magnetización en sistemas de Ising.

\subsection{GNN para Sistemas de Espines}

\citet{kochkov2021} emplearon GNN como manifold variacional para estados fundamentales de Hamiltonianos de Heisenberg, demostrando que la arquitectura respeta simetrías físicas por construcción. \citet{huang2022} proporcionaron la fundamentación teórica más rigurosa, demostrando que el aprendizaje automático clásico puede predecir eficientemente propiedades del estado fundamental de Hamiltonianos con gap, generalizando a nuevas instancias dentro de la misma fase.


\section{Justificación de Sistemas de Espines frente a Química Cuántica}

Una decisión arquitectónica fundamental de este trabajo es la orientación hacia sistemas de espines, en lugar de química cuántica molecular. Para sistemas de espines existe un isomorfismo natural entre espín-1/2 y qubit, con interacciones $Z_i Z_j$ que se traducen directamente en compuertas nativas. Para el TFIM con $N$ sitios y HVA $p = 2$, el circuito requiere $2(N-1)$ compuertas de dos qubits ---lineal en el tamaño del sistema. En contraste, la transformación de Jordan-Wigner para fermiones mapea operadores locales a cadenas de Pauli de longitud $O(N)$, incrementando la profundidad a $O(N^4)$ para ansätze tipo UCCSD \citep{tilly2022} ---superando ampliamente el límite de truncamiento $O(\log n)$ \citep{mele2026}. Además, la clasificación de fases cuánticas es una pregunta cualitativa con criterio relajado ($\Delta E/\text{gap} < 5\%$), mientras que la química cuántica exige precisión química ($\Delta E < 1.6 \times 10^{-3}$ Hartree $\approx 1$ kcal/mol) \citep{tilly2022}, inalcanzable con circuitos superficiales en hardware ruidoso.

\section{Síntesis y Brecha Identificada}

El análisis del estado del arte revela:

\begin{enumerate}
    \item El truncamiento por ruido \citep{mele2026} impone circuitos superficiales como \emph{requisito} para hardware real, motivando el uso de HVA de profundidad moderada.
    \item Las funciones de costo locales garantizan gradientes detectables \citep{cerezo2021}.
    \item La transferibilidad de parámetros \citep{mele2022, puig2025} y la capacidad de las GNN para predecir propiedades de estados fundamentales \citep{huang2022, zhang2025} fundamentan la viabilidad del predictor MPNN.
    \item La predicción de parámetros variacionales mediante ML constituye un paradigma consolidado con al menos siete trabajos independientes \citep{miao2024, zhang2025, zou2026, yang2024pvls, zhang2025vqezy, huang2025palqo, karim2025}, demostrando que la técnica individual no es novedosa per se.
    \item La mitigación de errores (PEA-ZNE + DD + TREX) ha sido validada experimentalmente en hardware IBM a escalas comparables \citep{kim2023, kiiamov2026, sharma2026}, estableciendo la viabilidad del despliegue futuro.
\end{enumerate}

La brecha identificada no reside en la técnica de predicción ---cuya viabilidad está demostrada por la literatura citada--- sino en la ausencia de una \emph{validación sistemática integrada} que responda preguntas prácticas: ¿para qué topologías funciona? ¿Hasta qué profundidad de circuito? ¿Para qué familias de Hamiltonianos? ¿Cuáles son los límites físicos del ansatz? Los trabajos existentes abordan componentes individuales ---predicción de parámetros \citep{miao2024, zhang2025, huang2025palqo}, diseño de ansatz \citep{wiersema2020, tripathi2026}, warm-start \citep{puig2025}, datasets \citep{zhang2025vqezy}--- pero la validación conjunta sobre 5 topologías distintas, múltiples profundidades de circuito ($p = 1$--$4$), variantes del Hamiltoniano, y la documentación formal de las fronteras de expresividad donde el enfoque deja de funcionar, no se ha reportado en la literatura revisada.

El presente trabajo contribuye con esta \emph{integración y validación sistemática}: un pipeline reproducible donde cada componente está implementado con técnicas del estado del arte, y cuyo valor reside en (i) la demostración de que la combinación funciona end-to-end con aceleraciones de $29$--$500\times$, (ii) la documentación cuantitativa de los límites de aplicabilidad (modelo, topología, profundidad, escala), y (iii) el establecimiento de una metodología de benchmark para futuros trabajos en el área.


%---------------------------
% CAPÍTULO 3: DESARROLLO DEL TRABAJO
%---------------------------

\chapter{Desarrollo del Trabajo}

Este capítulo describe la metodología completa del pipeline híbrido propuesto, incluyendo la arquitectura del software, los algoritmos de cada fase, las estrategias de validación y los criterios de éxito. Se detalla la implementación concreta que permite reproducir todos los resultados presentados en el Capítulo~4.

\section{Objetivo General}

Demostrar que la integración de predicción de parámetros variacionales mediante una Red Neuronal de Grafos de Paso de Mensajes (MPNN) con un Ansatz Variacional Hamiltoniano (HVA) de profundidad acotada ($p \leq 4$) y warm-start descendente, implementada como un pipeline híbrido clásico-cuántico unificado, permite reducir el costo cuántico de la clasificación de fases del Modelo de Ising en Campo Transversal entre $29\times$ y $500\times$ respecto a VQE convencional con inicialización aleatoria, manteniendo un error relativo $\Delta E/\text{gap} < 5\%$ dentro del régimen operativo válido, y documentar formalmente los límites de dicho régimen ---incluyendo la frontera de expresividad del ansatz, las condiciones de fallo por topología y modelo, y la ley de escalado del campo mínimo operable--- validado mediante simulación ideal sobre 5 topologías de red, tamaños de sistema $N = 4$--$20$, y extensiones del Hamiltoniano.

La contribución principal no reside en las técnicas individuales ---que fueron propuestas en la literatura por \citet{miao2024}, \citet{zhang2025}, \citet{wiersema2020} y \citet{puig2025}--- sino en su integración sistemática en un pipeline reproducible, la validación exhaustiva de su dominio de aplicabilidad ($400+$ ejecuciones, 5 topologías, 3 modelos), y la documentación rigurosa de los límites físicos y computacionales que determinan dónde el enfoque es viable y dónde fracasa.

\section{Objetivos Específicos}

\begin{enumerate}
    \item \textbf{OE1:} Analizar el estado del arte en algoritmos cuánticos variacionales y estrategias de inicialización basadas en aprendizaje automático, identificando las barreras teóricas (barren plateaus, truncamiento por ruido, costo de evaluación) y las condiciones bajo las cuales la combinación de circuitos superficiales, costos locales y warm-start las sortea simultáneamente \citep{mele2026, cerezo2021, wiersema2020}.

    \item \textbf{OE2:} Desarrollar un módulo de generación de ground truth parametrizable por topología y tamaño de sistema, mediante diagonalización exacta ($N \leq 14$) y DMRG ($N > 14$) \citep{hauschild2018}, que produzca datasets verificados para entrenamiento y validación del predictor.

    \item \textbf{OE3:} Implementar un pipeline VQE con warm-start descendente sobre HVA ---siguiendo la estrategia de transferibilidad de parámetros formalizada por \citet{mele2022} y \citet{puig2025}--- alcanzando fidelidades $\geq 93\%$ en al menos el 90\% de los puntos del barrido dentro del régimen operativo válido.

    \item \textbf{OE4:} Diseñar y entrenar un predictor MPNN basado en GINConv \citep{xu2019} ---siguiendo el paradigma de predicción de parámetros propuesto por \citet{miao2024} y extendido por \citet{zhang2025}--- que prediga parámetros óptimos con $\Delta E/\text{gap} < 5\%$ para múltiples topologías y tamaños de sistema, eliminando el bucle iterativo de optimización cuántica.

    \item \textbf{OE5:} Validar sistemáticamente el pipeline sobre 5 topologías (cadena 1D, escalera, triangular, cuadrada, heavy-hex) y tamaños $N = 4$ a $N = 20$, documentando reproducibilidad (multi-semilla), límites de expresividad ($h_{\min}$ por configuración), y condiciones de fallo, estableciendo una metodología de evaluación con criterios cuantitativos de aceptación/rechazo.

    \item \textbf{OE6:} Evaluar la extensibilidad del framework a variantes del Hamiltoniano de Ising (campo longitudinal, frustración $J_1$-$J_2$) y modelos de mayor complejidad (Heisenberg), documentando formalmente tanto los éxitos como los límites fundamentales de expresividad ---incluyendo el presupuesto de puertas CX, la incompatibilidad del estado inicial, y las barreras de frustración \citep{huang2026frustration}--- para establecer el dominio de aplicabilidad del enfoque.
\end{enumerate}

\section{Arquitectura del Pipeline}

El pipeline se organiza en tres fases secuenciales, cada una produciendo artefactos verificables que alimentan la siguiente. El diseño modular garantiza que cada componente puede ser reemplazado o escalado independientemente.

\subsection{Fase 1: Generación del Ground Truth Clásico}

\textbf{Objetivo:} Producir el dataset de referencia para entrenamiento y validación.

\textbf{Instrumentos:}
\begin{itemize}
    \item Diagonalización exacta (\texttt{numpy.linalg.eigh}) para $N \leq 14$.
    \item DMRG (TeNPy \citep{hauschild2018}) para $N > 14$, con $\chi = 64$--$200$ (exacto para HVA $p \leq 2$ en 1D, validado con diferencia $< 10^{-14}$).
    \item Hamiltoniano construido como \texttt{SparsePauliOp} (Qiskit 2.x).
    \item Generador de topologías: cadena 1D, escalera, triangular, cuadrada, heavy-hex.
\end{itemize}

\textbf{Procedimiento:}
\begin{enumerate}
    \item Definir topología ($N$ sitios, tipo de conectividad, condiciones de contorno).
    \item Construir el Hamiltoniano para un barrido no uniforme de $h$: $\Delta h = 0.1$ lejos del punto crítico y $\Delta h = 0.05$ en la región $h \in [0.8, 1.4]$, totalizando 27 puntos para $N \leq 10$.
    \item Para cada $h$: calcular estado fundamental $|\psi_0\rangle$, energía $E_0$, gap espectral $\Delta = E_1 - E_0$, y observables locales $\langle X_i \rangle$, $\langle Z_i Z_{i+1} \rangle$.
\end{enumerate}

\textbf{Criterio de éxito:} Dataset completo con $\geq 27$ puntos, energías verificadas (error $< 10^{-12}$ para diag. exacta, $< 10^{-8}$ para DMRG).

\subsection{Fase 2: Optimización VQE con Warm-Start}

\textbf{Objetivo:} Obtener los parámetros óptimos $\theta^*(h)$ del HVA para cada punto del barrido, generando el dataset de entrenamiento para la MPNN.

\textbf{Instrumentos:}
\begin{itemize}
    \item Circuitos HVA de profundidad $p = 1$--$4$ ($2p$ parámetros, independiente de $N$).

    \item \textbf{Evaluación del circuito:} \texttt{StatevectorEstimator} (Qiskit Primitives V2) para $N \leq 22$, donde el vector de estado tiene $2^N$ amplitudes ($\sim 64$\,MB a $N = 22$); simulación MPS (Aer, dimensión de enlace $\chi = 64$) para $N > 22$. El simulador MPS representa el estado cuántico como un producto tensorial de matrices con dimensión de enlace acotada, escalando como $O(N \cdot \chi^3)$ en lugar de $O(2^N)$ \citep{hauschild2018}. En modo determinista, la energía se calcula exactamente mediante contracción del MPS ($|\Delta E| < 10^{-14}$ respecto a statevector, validado experimentalmente). La dimensión $\chi = 64$ es exacta para circuitos HVA de profundidad moderada en 1D, ya que el entrelazamiento generado es bajo (verificado con diferencia $< 10^{-14}$ frente a statevector a $N = 22$).

    \item \textbf{Optimizador:} L-BFGS-B ($f_{\text{tol}} = 10^{-14}$) para simulación noiseless, que explota la evaluación exacta del gradiente \citep{tilly2022}; COBYLA para simulación estocástica (MPS con shots), que opera sin gradientes y es robusto al ruido de muestreo \citep{singh2025}.

    \item \textbf{Estrategia multi-restart:} Para cada punto $h$ del barrido, se ejecutan 3--7 reinicios del optimizador (según la escala del sistema: 3 reinicios para $N \leq 10$, 7 para $N = 20$ donde se detectan 2--3 mínimos locales). El primer reinicio hereda $\theta^*$ del punto anterior (warm-start); los restantes parten de perturbaciones gaussianas del warm-start ($\sigma = 0.1$). Se selecciona la solución de menor energía entre todos los reinicios.
\end{itemize}

\textbf{Procedimiento:}
\begin{enumerate}
    \item Inicializar $\theta_0 = \text{uniform}(-0.01, 0.01)$ para $h_{\max}$ (el estado $|+\rangle^{\otimes N}$ es casi exacto en la fase paramagnética profunda).
    \item Barrido descendente ($h = A \to B$, donde A > B): cada punto hereda $\theta^*$ del anterior (warm-start), siguiendo la transferibilidad de parámetros del HVA \citep{mele2022, puig2025}.
    \item Multi-restart en cada punto: warm-start heredado + perturbaciones gaussianas.
    \item Seleccionar la solución de menor energía entre reinicios; calcular fidelidad $|\langle \psi(\theta^*) | \psi_0 \rangle|^2$.
    \item Filtrar puntos con fidelidad $< 0.93$ (excluidos del entrenamiento MPNN, típicamente en la región crítica).
\end{enumerate}

\textbf{Criterio de éxito:} Fidelidad media $\geq 0.95$; $\geq 90\%$ de los puntos con fidelidad $\geq 0.93$ dentro del régimen válido; paisaje $\theta(h)$ suave ($\max_i \|\theta(h_i) - \theta(h_{i-1})\|_\infty < 1.0$).

\subsection{Fase 3: Entrenamiento del Predictor MPNN}

\textbf{Objetivo:} Entrenar una red neuronal de grafos que prediga $\theta^*$ directamente desde la representación en grafo del Hamiltoniano, eliminando la optimización VQE en tiempo de inferencia.

\textbf{Arquitectura:}
\begin{itemize}
    \item 3 capas GINConv \citep{xu2019} + \texttt{global\_mean\_pool} + MLP de salida. La agregación \texttt{global\_mean\_pool} promedia las representaciones de todos los nodos para producir un vector de tamaño fijo independiente de $N$.
    \item Framework: PyTorch Geometric \citep{fey2019}.
    \item Representación: nodos = qubits (atributo: $h_i$), aristas = interacciones ($J_{ij}$).
    \item Hiperparámetros: hidden\_dim $= 64$--$128$ (según $N$); epochs $= 6000$; lr $= 10^{-3}$; dropout $= 0.1$ (desactivación aleatoria del 10\% de neuronas durante entrenamiento para prevenir sobreajuste \citep{gal2016}); patience $= 300$--$500$.
\end{itemize}

\textbf{Procedimiento:}
\begin{enumerate}
    \item Construir dataset de grafos a partir de los pares $(h, \theta^*(h))$ filtrados por fidelidad.
    \item Dividir en entrenamiento (80\%) y validación (20\%).
    \item Entrenar minimizando MSE entre $\theta_{\text{pred}}$ y $\theta^*$.
    \item Validación física: evaluar la energía $E(\theta_{\text{pred}})$ y verificar $\Delta E/\text{gap} < 5\%$.
    \item Análisis de gradientes de pesos \citep{hernandes2025}: detectar picos en $\|\partial W / \partial h\|$ como indicador de transición de fase.
\end{enumerate}

\textbf{Criterio de éxito:} MSE $< 5 \times 10^{-3}$; $\Delta E/\text{gap} < 5\%$ en el régimen válido; detección de pico de gradiente en $h \approx 1.0$--$1.2$.

\textbf{Justificación de GINConv sobre alternativas:} La elección de GINConv se fundamenta en: (1) GIN es maximalmente expresiva entre las MPNNs \citep{xu2019}; (2) para redes con acoplamientos uniformes ($J$ constante), los mecanismos de atención (GATConv) añaden complejidad sin ganancia informativa; (3) \citet{meng2025} confirmaron independientemente la superioridad de GNN sobre CNN (+36\%) para predicción de propiedades de circuitos cuánticos.

\textbf{Variante UnifiedMPNN (Bond-Resolved, Multi-N):} Para habilitar la generalización a tamaños de sistema no vistos (cross-N), se implementa una variante denominada \emph{UnifiedMPNN} que opera sobre la parametrización bond-resolved del HVA ---donde cada enlace $(i,j)$ posee su propio ángulo $\theta_{ij}$--- en lugar de parámetros globales compartidos. La arquitectura difiere de la MPNN estándar en tres aspectos: (1) el target de predicción son los ángulos por enlace (dimensión variable con $N$) en lugar de un vector fijo de $2p$ valores; (2) se emplea BatchNorm en las capas ocultas para estabilizar el entrenamiento con grafos de tamaño heterogéneo; y (3) el entrenamiento es conjunto sobre múltiples valores de $N$, enseñando a la red la relación entre estructura local de cada enlace y su parámetro óptimo. Esta arquitectura aprende una función local ($\text{vecindario}(i,j), h \mapsto \theta_{ij}$) que generaliza naturalmente a grafos de tamaño arbitrario, superando la limitación del \texttt{global\_mean\_pool} que produce embeddings incompatibles entre tamaños distintos.

\subsection{Fase 3b: Evaluación de la Predicción (Deploy)}

\textbf{Objetivo:} Verificar que los parámetros predichos por la MPNN producen clasificación correcta de fases al ejecutar el circuito HVA en simulación ideal.

\textbf{Criterio de éxito:}
\begin{enumerate}
    \item $\Delta E/\text{gap} < 5\%$ (métrica primaria).
    \item $\Delta E$ absoluto comparable al obtenido por VQE directo.
    \item Fidelidad $\geq 0.95$ respecto al estado fundamental exacto.
    \item Clasificación correcta de fase (paramagnética vs. ferromagnética).
\end{enumerate}

\section{Extensibilidad a Variantes del Hamiltoniano de Ising}

El framework se diseña con un sistema de registro de modelos (\texttt{ModelRegistry}) que permite configurar nuevos Hamiltonianos sin modificar el pipeline:

\begin{itemize}
    \item \textbf{TFIM estándar} (Ec.~\ref{eq:tfim}): Sistema de validación principal. 10 CZ gates a $N=6$, $p=1$.
    \item \textbf{TFIM + campo longitudinal} (Ec.~\ref{eq:tfim_long}): El término $g \cdot Z_i$ se implementa con compuertas $R_Z$ (un qubit), sin incrementar el conteo de compuertas de dos qubits. Permite validar la extensibilidad sin costo hardware adicional.
    \item \textbf{TFIM frustrado $J_1$-$J_2$} (Ec.~\ref{eq:tfim_frustrated}): Los acoplamientos a segundos vecinos añaden compuertas CZ adicionales (27 CZ a $N=6$), limitando la viabilidad de ZNE para $N \geq 6$.
    \item \textbf{Heisenberg XXZ:} Evaluado como caso límite negativo para documentar las fronteras de expresividad del HVA $p=2$.
\end{itemize}

La regla de extensibilidad se fundamenta en la estructura de compilación del circuito HVA: los términos del Hamiltoniano que son producto tensor de operadores de un qubit ($X_i$, $Z_i$) se implementan con rotaciones de un qubit ($R_X$, $R_Z$), sin incrementar el conteo de compuertas de dos qubits \citep{wiersema2020}. Modelos que requieren interacciones de dos cuerpos adicionales (XX, YY) introducen compuertas CZ/CX por cada término, excediendo el presupuesto de mitigación para profundidades moderadas (validado en 46 runs Heisenberg con 0\% PassRate, véase Sección~\ref{sec:extensibilidad}).

\section{Marco de Validación Sistemática}

La validación del pipeline se estructura en tres niveles de comprobación progresiva:

\textbf{Nivel 1 --- Pruebas unitarias y de integración:} Suite de 72 tests automatizados (pytest, §\ref{sec:codigo}) que verifican cada módulo individualmente y la integración del pipeline completo. Incluye pruebas de humo (\texttt{smoke\_test}: $N=4$, $p=1$, $<30$s) ejecutadas en integración continua.

\textbf{Nivel 2 --- Validación estadística multi-semilla:} Cada configuración se ejecuta con mínimo 3 semillas aleatorias (42, 43, 44) para evaluar reproducibilidad. Se reportan media $\pm$ desviación estándar de $\Delta E/\text{gap}$.

\textbf{Nivel 3 --- Validación cruzada por topología y escala:} El pipeline se ejecuta sobre todas las combinaciones de topología $\times$ tamaño $\times$ profundidad disponibles, produciendo una matriz de resultados que permite identificar patrones de éxito y fallo.

\subsection{Diagnóstico Automatizado de Fallos}

Se implementa un sistema de diagnóstico que clasifica automáticamente las ejecuciones fallidas en categorías de causa raíz:

\begin{itemize}
    \item \textbf{CHAIN\_BREAK} ($\theta_{\text{smooth}} > 1.0$): Discontinuidad en la cadena de warm-start, detectable en Fase 2.
    \item \textbf{MPNN\_OVERFIT} (gen\_gap $> 0.01$): Sobreajuste del predictor, detectable en Fase 3.
    \item \textbf{BOUNDARY\_EFFECT}: Punto de test demasiado cercano al límite del régimen válido.
    \item \textbf{OUTSIDE\_REGIME}: Punto de test por debajo del régimen válido.
    \item \textbf{VQE\_DIVERGENCE}: No convergencia del optimizador.
\end{itemize}

El pipeline de early-stopping (alertar si $\theta_{\text{smooth}} > 1.0$, abortar si gen\_gap $> 0.01$) previene el 69\% de los fallos sin rechazar ninguna configuración exitosa (véase Figura~\ref{fig:failure_treemap} en Capítulo~5).

\section{Herramientas y Entorno de Desarrollo}

\begin{itemize}
    \item \textbf{Lenguaje:} Python 3.11+.
    \item \textbf{Computación cuántica:} Qiskit 2.x (Primitives V2, \texttt{SparsePauliOp}, pass managers).
    \item \textbf{Hardware cuántico:} IBM Quantum Platform (procesadores Torino/Heron, 133+ qubits).
    \item \textbf{Redes tensoriales:} TeNPy \citep{hauschild2018} para DMRG.
    \item \textbf{Machine Learning:} PyTorch + PyTorch Geometric \citep{fey2019}.
    \item \textbf{Optimización:} SciPy (L-BFGS-B, COBYLA).
    \item \textbf{Control de calidad:} 72 tests unitarios, linter (ruff), CI/CD (GitHub Actions), hooks pre-commit.
    \item \textbf{Arquitectura modular:} Paquete instalable \texttt{qmbp\_simulation} con 12 módulos independientes bajo DAG estricto de dependencias.
\end{itemize}


%---------------------------
% CAPÍTULO 4: DESCRIPCIÓN DETALLADA E IDENTIFICACIÓN DE REQUISITOS
%---------------------------

\chapter{Descripción Detallada e Identificación de Requisitos}

Este capítulo detalla el problema computacional abordado, su contexto dentro de la simulación cuántica de muchos cuerpos, y los requisitos funcionales que la solución debe satisfacer. Se describe además el trabajo previo realizado para delimitar el alcance del pipeline propuesto.

\section{Contexto del Problema}

La caracterización de fases cuánticas en hardware NISQ enfrenta un conflicto fundamental entre tres restricciones simultáneas \citep{tilly2022, mele2026}:

\begin{enumerate}
    \item \textbf{Restricción de profundidad:} El ruido no unital limita la profundidad útil del circuito a $d_{\text{eff}} = O(\log(n/\varepsilon)/\gamma)$ capas \citep{mele2026}. Cada capa adicional incrementa el presupuesto de puertas de dos qubits y el costo exponencial de mitigación \citep{tsubouchi2023}, imponiendo la necesidad de mantener $p$ acotado.

    \item \textbf{Restricción de optimización:} VQE convencional requiere $\sim 500$--$1000$ evaluaciones por punto del diagrama de fases \citep{tilly2022}. Para un barrido de 27 puntos en $h$, esto implica $\sim 15{,}000$--$27{,}000$ ejecuciones de circuito ---prohibitivo en QPU con costo por shot y limitaciones de cola.

    \item \textbf{Restricción de expresividad:} El HVA con profundidad moderada no puede representar estados fundamentales en la región crítica ($h \approx h_c$) para ningún tamaño de sistema \citep{tripathi2026, sumeet2025}, limitando el régimen operativo a $h \geq h_{\min}(N)$. La profundidad necesaria para convergencia exacta escala como $p \approx N/2$ \citep{sumeet2025}.
\end{enumerate}

El problema concreto es: ¿puede eliminarse la optimización iterativa en hardware cuántico sin sacrificar la precisión en la clasificación de fases, operando dentro del régimen donde el HVA superficial es expresivo?

\section{Trabajo Previo}

El desarrollo del pipeline se realizó en fases iterativas documentadas:

\begin{itemize}
    \item \textbf{Versiones 1--5:} Exploración de ansätze, funciones de costo y estrategias de inicialización, definición de principios de diseño como la necesidad de funciones de costo puras (energía) para pipelines ML-VQE \citep{miao2024}.

    \item \textbf{Versión 6:} Estabilización del pipeline con función de costo de energía pura, filtro de fidelidad $\geq 0.93$, y validación para $N = 6$ y $N = 10$ (40+ experimentos por escala).

    \item \textbf{Versiones 7--9:} Extensión sistemática a múltiples topologías, Heisenberg como caso negativo, confirmando límites de expresividad para circuitos reducidos.

    \item \textbf{S-series:} Experimentos focalizados (S1--S8) para eficiencia de datos, paisaje energético, incertidumbre y escalamiento.
\end{itemize}


%---------------------------
% CAPÍTULO 5: EVALUACIÓN
%---------------------------

\chapter{Evaluación}

Este capítulo presenta la evaluación completa del pipeline GNN-HVA mediante simulación ideal (sin ruido), evaluando la ventaja de la solución propuesta y su aplicabilidad para resolver el problema de clasificación de fases cuánticas. Todos los resultados se obtienen con el backend \texttt{StatevectorEstimator} (Qiskit Primitives V2), que calcula valores esperados de forma exacta sin shots ni modelo de ruido, estableciendo el techo teórico del pipeline. Se organizan los resultados en cinco ejes: (1) validación por topología a escala fija, (2) validación de escalamiento con profundidad y tamaño del sistema, (3) extensibilidad a variantes del Hamiltoniano de Ising, (4) comparativa con la literatura, y (5) análisis de límites físicos del ansatz. La totalidad de los resultados numéricos se basa en 400+ ejecuciones del pipeline sobre 5 topologías, profundidades $p = 1$--$4$, y tamaños de sistema desde $N = 4$ hasta $N = 20$.

\section{Definición de Métricas}

Se definen formalmente las métricas empleadas a lo largo de este capítulo:

\begin{itemize}
    \item \textbf{$\Delta E/\text{gap}$} (Error relativo al gap espectral): Métrica primaria de éxito, definida como:
    \begin{equation}
    \frac{\Delta E}{\text{gap}} = \frac{|E_{\text{pred}} - E_{\text{exact}}|}{E_1 - E_0}
    \end{equation}
    donde $E_{\text{pred}} = \langle \psi(\theta_{\text{pred}}) | H | \psi(\theta_{\text{pred}}) \rangle$ es la energía obtenida con los parámetros predichos por la MPNN, $E_{\text{exact}} = E_0$ es la energía del estado fundamental (calculada por diagonalización exacta o DMRG), y $E_1 - E_0$ es el gap espectral. Normalizar por el gap permite comparar entre distintos valores de $h$ y tamaños $N$: cerca del punto crítico el gap se cierra y el error relativo aumenta, reflejando la mayor dificultad física. El umbral de éxito $\Delta E/\text{gap} < 5\%$ garantiza clasificación correcta de fase.

    \item \textbf{$\Delta E$} (Error energético absoluto): Diferencia absoluta entre la energía predicha y la exacta:
    \begin{equation}
    \Delta E = |E_{\text{pred}} - E_{\text{exact}}|
    \end{equation}
    Mide cuán lejos está la predicción del estado fundamental en unidades de energía, sin normalizar por el gap. Es complementaria a $\Delta E/\text{gap}$: un valor $\Delta E$ pequeño pero $\Delta E/\text{gap}$ grande indica proximidad al punto crítico donde el gap se cierra, no un fallo del pipeline.

    \item \textbf{Fidelidad} ($F$): Overlap cuadrático entre el estado preparado y el estado fundamental exacto:
    \begin{equation}
    F = |\langle \psi(\theta_{\text{pred}}) | \psi_{\text{exact}} \rangle|^2
    \end{equation}
    Mide la calidad del estado cuántico predicho independientemente de la energía. Dos estados pueden tener energías similares pero fidelidad baja si pertenecen a ramas distintas del espacio de Hilbert. Disponible únicamente en simulación ideal (requiere acceso al vector de estado completo).

    \item \textbf{PassRate} (Tasa de aprobación): Fracción de puntos $h$ evaluados que satisfacen $\Delta E/\text{gap} < 5\%$. Un pipeline exitoso alcanza PassRate $\geq 90\%$ dentro del régimen operativo válido.

    \item \textbf{Speedup} (Aceleración respecto a VQE random): Ratio entre las iteraciones requeridas por VQE con inicialización aleatoria y las evaluaciones necesarias con el pipeline. Se calcula como el número de evaluaciones de circuito evitadas gracias a la predicción directa de la MPNN.

\end{itemize}

\section{Validación del Pipeline Base ($N = 6$, $N = 10$)}

\subsection{Resultados para $N = 6$ y $N = 10$ (Cadena 1D, $p = 2$)}

La validación inicial se realiza en puntos de test cercanos al punto crítico ($h_c = 1.0$), donde el gap espectral es pequeño y la clasificación es difícil. Esto constituye un test exigente del pipeline:
Se observa que $\Delta E$ (error absoluto) es comparable entre puntos ($\sim$0.015--0.028), pero $\Delta E/\text{gap}$ varía significativamente porque el gap se cierra al acercarse a $h_c$. A $h = 1.25$ ($N = 6$), el gap es solo 0.46 y el pipeline opera al límite de su criterio de éxito (3.77\%). Esto confirma que la frontera de expresividad del HVA determina el régimen válido.

\subsection{Resultados para $N = 10$, $p = 1$}

Con profundidad mínima ($p = 1$, 2 parámetros), el pipeline funciona en las 4 topologías. Nótese que los $h_{\text{test}}$ están en la fase paramagnética profunda ($h \gg h_c$), donde el gap es grande y el paisaje es trivial --- estos puntos representan la condición operativa más favorable:

\begin{table}[h]
\centering
\caption{Pipeline $p = 1$, $N = 10$, 4 topologías (media $\pm$ std sobre 3 semillas). Los puntos de test están lejos de $h_c$: gap grande, predicción más sencilla.}
\label{tab:results_p1}
\begin{tabular}{lccc}
\toprule
\textbf{Topología} & \textbf{$h_{\text{test}}$} & \textbf{Mediana $\Delta E/\text{gap}$} & \textbf{Pasan} \\
\midrule
Cadena 1D & 2.75 & 4.10\% & 3/3 \\
Escalera & 3.25 & 2.90\% $\pm$ 0.40\% & 3/3 \\
Triangular & 4.25 & 3.30\% & 3/3 \\
Heavy-Hex & 3.25 & 0.56\% $\pm$ 0.03\% & 3/3 \\
\bottomrule
\end{tabular}
\end{table}

La comparación entre ambas tablas es instructiva: a $p = 1$ con $h = 2.75$--$4.25$ (gap $\approx 3$--$7$) se obtiene $\Delta E/\text{gap} \approx 3\%$, mientras que a $p = 2$ con $h = 1.25$--$1.50$ (gap $\approx 0.5$--$1.2$) se obtiene un resultado similar (2--5\%). El error absoluto $\Delta E$ es el factor estable; la dificultad aparente surge del gap como denominador.



\section{Validación Cross-Topology para $N = 10$}

El pipeline se evaluó sobre cinco topologías para $N = 10$ con profundidades $p = 1$--$4$ (5 topologías $\times$ 4 profundidades $\times$ 3--4 semillas = 79 ejecuciones). El rango de evaluación $h \in [1.0, 5.0]$ (Tabla~\ref{tab:cross_topo_depth}) abarca desde la proximidad al punto crítico (gap pequeño, predicción difícil) hasta la fase paramagnética profunda (gap grande, predicción trivial). La Tabla~\ref{tab:cross_topo} reporta los mejores resultados restringidos al régimen válido ($h \geq 1.3$).

\begin{table}[h]
\centering
\caption{Mejores configuraciones por topología ($N = 10$, TFIM, $h \in [1.3, 3.0]$, 40 puntos).}
\label{tab:cross_topo}
\begin{tabular}{lcccc}
\toprule
\textbf{Topología} & \textbf{Mejor $p$} & \textbf{PassRate} & \textbf{$\bar{F}_{\text{VQE}}$} & \textbf{Speedup} \\
\midrule
Cadena 1D & 3 & \textbf{95\%} (37/39) & 0.997 & 44$\times$ \\
Heavy-Hex & 3 & \textbf{90\%} (35/39) & 0.995 & 82$\times$ \\
Escalera & 4 & 79\% (31/39) & 0.945 & 182$\times$ \\
Cuadrada & 4 & 79\% (31/39) & 0.953 & 61$\times$ \\
Triangular & 4 & 56\% (22/39) & 0.966 & 43$\times$ \\
\bottomrule
\end{tabular}
\end{table}

\begin{table}[h]
\centering
\caption{Progresión del PassRate ($\Delta E/\text{gap} < 5\%$) con la profundidad del circuito por topología ($N = 10$, TFIM, $h \geq 1.3$).}
\label{tab:cross_topo_depth}
\begin{tabular}{lcccc}
\toprule
\textbf{Topología} & \textbf{$p = 1$} & \textbf{$p = 2$} & \textbf{$p = 3$} & \textbf{$p = 4$} \\
\midrule
Cadena 1D & 78\% & 92\% & \textbf{100\%} & \textbf{100\%} \\
Heavy-Hex & 81\% & 92\% & 97\% & \textbf{100\%} \\
Escalera & 58\% & 75\% & 81\% & 86\% \\
Cuadrada & 53\% & 69\% & 81\% & 86\% \\
Triangular & 28\% & 39\% & 53\% & 61\% \\
\bottomrule
\end{tabular}
\end{table}

\textbf{Hallazgo 1: La topología domina sobre la profundidad del circuito.} Para topologías 2D (escalera, cuadrada, triangular), incrementar $p$ de 1 a 4 mejora el PassRate en +28--33 puntos porcentuales, mientras que cadena 1D y heavy-hex alcanzan 100\% desde $p=3$ (Tabla~\ref{tab:cross_topo_depth}). La conectividad del grafo ---no la profundidad del ansatz--- es el factor limitante: incluso a $p=4$, escalera y cuadrada saturan en 86\% y triangular en 61\%.

\textbf{Hallazgo 2: Jerarquía de topologías invariante a $p$.} Se establece un ranking estable: cadena 1D $\approx$ heavy-hex $\gg$ escalera $\approx$ cuadrada $\gg$ triangular. La explicación es el número de compuertas CZ por capa del HVA: cadena 1D requiere $N-1 = 9$ CZ, heavy-hex $\sim$12, mientras que escalera/cuadrada/triangular requieren $15$+ CZ, creando paisajes de optimización con más mínimos locales.

\textbf{Hallazgo 3: Mejora monotónica del PassRate con $p$.} El PassRate crece consistentemente con la profundidad en todas las topologías cuando se evalúa en el régimen operativo válido ($h \geq 1.3$). El salto cualitativo ocurre en $p=3$, donde cadena 1D y heavy-hex superan 97\% (Tabla~\ref{tab:cross_topo_depth}).

\textbf{Hallazgo 4: La transferencia cross-topology falla.} La MPNN entrenada sobre datos de una topología \emph{no generaliza} a topologías distintas. Un modelo entrenado con datos de cadena 1D produce $\Delta E/\text{gap} \gg 5\%$ al evaluarse sobre escalera o triangular ($N = 10$, 3 semillas), confirmando que cada topología induce parámetros $\theta^*$ cualitativamente distintos para un mismo $h$.

\textbf{Hallazgo 5: $\Delta E$ absoluto discrimina entre fallo de expresividad y fallo de predicción.} La Tabla~\ref{tab:cross_topo_de} muestra el error energético absoluto y la fidelidad por topología (mejor configuración $p=3$--$4$). Para cadena 1D y heavy-hex, $\Delta E \approx 0.003$--$0.010$, indicando excelente proximidad al estado fundamental. Para topologías 2D, $\Delta E$ es un orden de magnitud mayor ($\sim$0.07--0.09) con fidelidad $F < 0.80$: el estado variacional no aproxima el ground state, lo que constituye un fallo de expresividad del ansatz \citep{tripathi2026}, no del predictor MPNN.

\begin{table}[h]
\centering
\caption{Error energético absoluto y fidelidad por topología ($N = 10$, TFIM, mejor configuración $p=3$--$4$, $h \geq 1.3$).}
\label{tab:cross_topo_de}
\begin{tabular}{lcccc}
\toprule
\textbf{Topología} & \textbf{$\Delta E$ (media)} & \textbf{$\Delta E$ (máx)} & \textbf{$F$ (media)} & \textbf{$F$ (mín)} \\
\midrule
Cadena 1D & 0.003 & 0.015 & 0.999 & 0.995 \\
Heavy-Hex & 0.010 & 0.044 & 0.997 & 0.984 \\
Escalera & 0.080 & 0.476 & 0.799 & 0.498 \\
Cuadrada & 0.089 & 0.240 & 0.776 & 0.499 \\
Triangular & 0.065 & 0.344 & 0.968 & 0.800 \\
\bottomrule
\end{tabular}
\end{table}

\textbf{Hallazgo 6: $\Delta E$ absoluto mejora con $p$ en todas las topologías.} Aunque $\Delta E/\text{gap}$ no satisface el umbral del 5\% para topologías 2D, el error energético absoluto sí decrece sistemáticamente al aumentar $p$ (Tabla~\ref{tab:delta_e_vs_p}). La mediana de $\Delta E$ se reduce 15--55$\times$ de $p=1$ a $p=4$. La discrepancia con $\Delta E/\text{gap}$ se explica por gaps menores en grafos de alta conectividad, que amplifican errores absolutos moderados. Para triangular, el error final (0.020) permanece un orden de magnitud superior al de cadena 1D (0.001), confirmando una limitación genuina del ansatz \citep{tripathi2026}.

\begin{table}[h]
\centering
\caption{Progresión de $\Delta E$ absoluto (mediana) con la profundidad $p$ por topología ($N = 10$, TFIM, $h \geq 1.3$). La mediana es robusta a outliers en la frontera del régimen válido.}
\label{tab:delta_e_vs_p}
\begin{tabular}{lccccc}
\toprule
\textbf{Topología} & \textbf{$p = 1$} & \textbf{$p = 2$} & \textbf{$p = 3$} & \textbf{$p = 4$} & \textbf{Mejora $p$1$\to$4} \\
\midrule
Cadena 1D   & 0.055 & 0.009 & 0.002 & 0.001 & 55$\times$ \\
Heavy-Hex   & 0.030 & 0.008 & 0.006 & 0.002 & 15$\times$ \\
Escalera    & 0.229 & 0.027 & 0.008 & 0.007 & 33$\times$ \\
Cuadrada    & 0.157 & 0.025 & 0.007 & 0.005 & 31$\times$ \\
Triangular  & 0.516 & 0.290 & 0.095 & 0.020 &  26$\times$ \\
\bottomrule
\end{tabular}
\end{table}

\noindent El salto entre $p=2$ y $p=3$ se explica por la suavidad del landscape: a $p \geq 3$, los parámetros óptimos $\theta^*(h)$ varían monotónicamente con $h$ (discontinuidades $< 0.05$ rad entre puntos consecutivos), permitiendo que la MPNN interpole con precisión. A $p \leq 2$, el landscape presenta saltos de $\pi$ (wraps periódicos) que la MPNN no puede capturar, resultando en predicciones con error dominado por la varianza del predictor.

\section{Validación de Escalamiento}

Se valida la escalabilidad del pipeline en dos dimensiones: (a) tamaño del sistema ($N = 6$--$20$ con $p = 3$--$4$, y $N = 6$--$200$ con $p = 1$), y (b) generalización a tamaños no vistos (zero-shot cross-N).

\subsection{Escalamiento con $N$ y $p$}

\textbf{Resultado principal:} Con $h$-range restringido al régimen válido ($h \geq 1.3$), el pipeline alcanza \textbf{100\% de tasa de aprobación} ($\Delta E/\text{gap} < 5\%$ en todos los puntos test) para $N = 10$, $16$ y $20$:

\begin{table}[h]
\centering
\caption{Mejores configuraciones por tamaño del sistema (cadena 1D, $h \geq 1.3$, pipeline completo). Backend: StatevectorEstimator ($N \leq 22$).}
\label{tab:scaling}
\begin{tabular}{rcccccc}
\toprule
\textbf{$N$} & \textbf{$p$} & \textbf{PassRate} & \textbf{$\Delta E$ (media)} & \textbf{$\bar{F}$} & \textbf{Speedup} \\
\midrule
8 & 4 & 97\% (28/29) & 0.006 & 0.999 & 63$\times$ \\
10 & 3 & \textbf{100\%} (24/24) & 0.003 & 0.999 & 29$\times$ \\
16 & 4 & \textbf{100\%} (34/34) & 0.006 & 0.998 & 64$\times$ \\
20 & 4 & \textbf{100\%} (39/39) & 0.005 & 0.995 & 58$\times$ \\
\bottomrule
\end{tabular}
\end{table}

El error absoluto $\Delta E$ se mantiene estable entre 0.003 y 0.006 independientemente de $N$, con fidelidad $\geq 0.995$ en todos los casos. Esto confirma que la calidad del estado predicho es uniformemente alta dentro del régimen válido, y que la mejora aparente de $\Delta E/\text{gap}$ con $N$ se debe al crecimiento del gap (denominador), no a una mejora del predictor.

Con $p = 1$ (2 parámetros, configuración mínima), el pipeline escala hasta $N = 250$ mediante simulación MPS ($\chi = 64$), con $\Delta E/\text{gap}$ que decrece de 2.70\% ($N = 10$) a 0.019\% ($N = 120$--$200$). Este descenso refleja que el gap crece con $h$ (y $h_{\min}$ crece con $N$), no una mejora intrínseca del predictor: el error absoluto $|\Delta E|$ se mantiene en $\sim$0.05--0.07 independientemente de $N$, y $|\Delta E|/N$ decrece de 0.0023 ($N = 30$) a 0.0009 ($N = 60$), indicando convergencia per-site. Con $p = 2$ (4 parámetros), la precisión mejora $70\times$ respecto a $p = 1$ en el mismo régimen ($\Delta E/\text{gap} = 0.005\%$ vs.\ 0.38\% a $N = 40$, backend MPS determinista, 3 semillas).

El pipeline end-to-end (Phase 1--3) se valida para $N = 30$--$60$ con la MPNN logrando $\Delta E/\text{gap} < 5\%$ en todos los puntos de deploy (interpolación), requiriendo solo 16--39 puntos de entrenamiento ---una eficiencia de datos notable que se atribuye a la suavidad de la función $h \mapsto \theta^*$ en el régimen paramagnético.

\textbf{Frontera del régimen válido ($h_{\min}$).} Para cada tamaño de sistema $N$ y profundidad $p$, existe un valor mínimo de $h$ por debajo del cual el HVA no puede representar el estado fundamental ---independientemente del optimizador o la inicialización. Este límite se origina en la capacidad de entrelazamiento finita del circuito: cerca del punto crítico ($h \to h_c = 1$), la entropía de entrelazamiento del ground state excede la que el ansatz de profundidad fija puede generar. \citet{tripathi2026} observaron independientemente este comportamiento para HVA en TFIM 1D/2D/3D, y \citet{sumeet2025} demostraron teóricamente que se requieren $p \approx N/2$ capas para convergencia exacta en el límite termodinámico.

La frontera se determinó empíricamente con barridos de alta resolución (10--20 puntos de $h$, 3 semillas, 5--7 restarts, backend MPS determinista $\chi = 64$). Los ajustes lineales para cadena 1D son:
\begin{equation}
h_{\min}(p{=}1) = 2.36 + 0.0073 \cdot N \quad (R^2 = 0.91, \; N = 20\text{--}250)
\label{eq:frontier_p1}
\end{equation}
\begin{equation}
h_{\min}(p{=}2) = 1.57 + 0.0050 \cdot N \quad (R^2 = 0.95, \; N = 20\text{--}120)
\label{eq:frontier_p2}
\end{equation}
Para $p \geq 3$, la frontera es cuasi-constante: $h_{\min}(p{=}3) \approx 1.6 \pm 0.1$ y $h_{\min}(p{=}4) \approx 1.4 \pm 0.1$, independientes de $N$. Este comportamiento es consistente con la \emph{area law}: en 1D, la entropía de entrelazamiento a $h/h_c$ fijo no depende del tamaño del sistema, y si el HVA tiene suficientes capas para capturar ese entrelazamiento, la frontera no se desplaza al escalar $N$. Para $p = 2$, la frontera baja a $h_{\min} \approx 1.7$--$1.9$ ($N = 20$--$40$), extendiendo el régimen operativo en $\sim$1.0 unidad de $h$ respecto a $p = 1$.

\begin{figure}[h]
\centering
\includegraphics[width=0.85\textwidth]{tesis-figures/fig_scaling_law_comprehensive.pdf}
\caption{Frontera del régimen válido (backend MPS determinista, 3 semillas, $\Delta E/\text{gap} = 5\%$). Para $p = 1$--$2$, $h_{\min}$ crece linealmente con pendiente $\sim$0.005--0.007 por qubit. Para $p \geq 3$, $h_{\min}$ se estabiliza (cuasi-constante), consistente con la area law de entrelazamiento en 1D.}
\label{fig:scaling_law}
\end{figure}

\subsection{Generalización Zero-Shot Cross-N}

La capacidad de la MPNN para generalizar se evalúa en dos escenarios: (a) interpolación intra-N (mismos $N$ del entrenamiento, valores de $h$ no vistos) y (b) generalización cross-N (tamaños de sistema no vistos).

\textbf{Intra-N (interpolación):} Entrenando con datos de un solo $N$ y evaluando en valores de $h$ no vistos, la MPNN logra 100\% PassRate para $N = 30$--$60$ con solo 16--39 puntos de entrenamiento ($p = 1$, cadena 1D, backend MPS $\chi = 64$). El error per-site $|\Delta E|/N$ decrece de 0.0023 ($N = 30$) a 0.0009 ($N = 60$), confirmando eficiencia de datos y escalado extensivo en modo intra-N.

\textbf{Cross-N con MPNN estándar (fallo):} Entrenando con datos de $N = 10, 20, 40$ y evaluando en $N = 60$ (zero-shot, $p = 1$, cadena 1D), la predicción con parametrización uniforme \textbf{falla} ($\Delta E/\text{gap} = 688\%$ media). La causa es que con 2 parámetros globales compartidos, $\theta^*(h)$ depende implícitamente de $N$ y el \texttt{global\_mean\_pool} produce embeddings incompatibles entre tamaños distintos.

\textbf{Cross-N con UnifiedMPNN bond-resolved (éxito):} Para superar esta limitación, se implementa una variante \emph{UnifiedMPNN} que emplea parametrización bond-resolved ---donde cada enlace $(i,j)$ recibe su propio ángulo $\theta_{ij}$--- combinada con BatchNorm y entrenamiento multi-N conjunto. Esta arquitectura aprende la relación entre la estructura local de cada enlace y su parámetro óptimo, generalizando naturalmente a grafos de tamaño arbitrario.

\begin{table}[h]
\centering
\caption{Validación cross-N con UnifiedMPNN bond-resolved ($p = 1$, TFIM, backend MPS $\chi = 64$). PassRate medido dentro del régimen válido ($h \geq h_{\min}$).}
\label{tab:cross_n_transfer}
\begin{tabular}{lccc}
\toprule
\textbf{Topología} & \textbf{$N$ entrenamiento} & \textbf{Total puntos} & \textbf{PassRate@5\%} \\
\midrule
Cadena 1D & 6,8,10,12,15,20 & 168 & 100\% \\
Heavy-Hex & 4,6,10,12,16 & 217 & 98\% \\
Escalera & 4,6,8,10,12,16,20 & 384 & 100\% \\
Cuadrada & 4,6,8,10,12,16,20 & 274 & 100\% \\
Triangular & 3,4,6,8,10,12 & 288 & 86\% \\
\bottomrule
\end{tabular}
\end{table}

El resultado confirma el claim principal: la GNN predice correctamente los ángulos del HVA dentro del régimen operativo válido de cada topología. La variante UnifiedMPNN resuelve la limitación cross-N del enfoque original, permitiendo entrenar un único modelo con datos de múltiples $N$ y obtener predicciones con $\Delta E/\text{gap} < 5\%$ consistentemente. La topología triangular (coordinación $z = 6$) muestra un PassRate menor (86\%) debido a que su régimen válido es más estrecho (la frontera $h_{\min}$ es más alta).

\textbf{Hallazgo técnico:} La canonicalización de simetría gauge ($\theta \mapsto \theta \bmod \pi$ + convención $\mathbb{Z}_2$) es un requisito para el entrenamiento multi-N. Sin ella, semillas distintas convergen a mínimos equivalentes pero con $\theta$ en periodos diferentes, contaminando los targets de la red.

\subsection{Escalado con el Tamaño del Sistema}
\label{subsec:large_n_extrap}

Para verificar que la calidad de predicción se mantiene al crecer $N$, se evalúa la UnifiedMPNN entrenada con datos de $N \leq 20$ sobre sistemas mayores ($N = 30$--$100$, backend MPS $\chi = 64$, ground truth DMRG). La evaluación se restringe al régimen válido ($h \geq h_{\min}$) de cada tamaño.

\begin{table}[h]
\centering
\caption{Predicción a large-N (cadena 1D, $p = 1$, UnifiedMPNN, backend MPS). Evaluación restringida al régimen válido.}
\label{tab:large_n_chain}
\begin{tabular}{rccccc}
\toprule
\textbf{$N$} & \textbf{$h$-range} & \textbf{Pts} & \textbf{$\Delta E/\text{gap}$ (media)} & \textbf{$|\Delta E|/N$} & \textbf{PassRate@5\%} \\
\midrule
10 & [3.5, 5.0] & 8 & 1.17\% & 8.1$\times 10^{-3}$ & 8/8 \\
20 & [3.5, 5.0] & 6 & 2.63\% & 9.1$\times 10^{-3}$ & 6/6 \\
30 & [2.5, 5.5] & 18 & 4.82\% & 9.4$\times 10^{-3}$ & 13/18 \\
40 & [2.5, 5.5] & 12 & 6.17\% & 8.6$\times 10^{-3}$ & 8/12 \\
\bottomrule
\end{tabular}
\end{table}

El error per-site $|\Delta E|/N$ permanece esencialmente constante ($\approx 9 \times 10^{-3}$) independientemente de $N$, confirmando escalado extensivo. Los puntos que no pasan el criterio del 5\% son aquellos cuyo valor de $h$ se encuentra cerca de la frontera $h_{\min}(N)$ ---no representan una degradación del predictor, sino un reflejo de que la frontera se desplaza con $N$ para $p = 1$ (Ec.~\ref{eq:frontier_p1}).

Este resultado valida que la GNN aprende correctamente el mapeo estructura-del-grafo $\to$ ángulos-óptimos, y que esta predicción se mantiene precisa al escalar el sistema siempre que se opere dentro del régimen donde el HVA es expresivo.

\subsection{Bond-Resolved HVA: Necesidad de la GNN en Alta Dimensionalidad}
\label{subsec:bond_resolved}

En el HVA estándar utilizado en este trabajo, todos los enlaces de un mismo tipo comparten un único parámetro variacional por capa ---por ejemplo, un solo ángulo $\theta_{ZZ}^{(l)}$ controla todas las compuertas $R_{ZZ}$ de la capa $l$. Esta parametrización \emph{uniforme} produce vectores $\theta$ compactos (2--4 parámetros para $p = 1$--$2$) y paisajes energéticos suaves, pero asume que la física es homogénea sobre todos los enlaces de la red.

Cuando la topología o el Hamiltoniano rompen esta homogeneidad ---por ejemplo, en redes bidimensionales irregulares, sistemas con desorden, o Hamiltonianos con acoplamientos espacialmente variables--- se requiere una parametrización \emph{bond-resolved}: cada enlace $(i,j)$ recibe su propio ángulo $\theta_{ij}^{(l)}$ independiente. Esto incrementa drásticamente la dimensionalidad del espacio de parámetros: para una cadena 1D con $N = 40$ y $p = 2$, se pasa de 4 parámetros (uniforme) a 79 parámetros (bond-resolved).

Se evalúa este régimen para determinar en qué punto la MPNN deja de ser una conveniencia y se convierte en una necesidad metodológica ---es decir, cuándo la interpolación clásica deja de ser viable como alternativa.

\textbf{Resultado:} Con 79 dimensiones ($N = 40$, $p = 2$, cadena 1D, bond-resolved, backend statevector, 3 semillas), la interpolación clásica es geométricamente impracticable. La GNN demuestra viabilidad intra-N con un factor de aceleración de $4414\times$ sobre inicialización aleatoria (6/6 PASS, $\Delta E/\text{gap} = 0.71\%$). Sin embargo, la generalización cross-N en 79D falla por insuficiencia de datos (45 puntos de entrenamiento para 494K parámetros de modelo). Este resultado define la frontera de investigación futura.

\subsection{Transferencia Directa de Parámetros VQE entre Tamaños: Irrelevante a Baja Dimensionalidad}

A diferencia de la predicción zero-shot de la MPNN (subsección anterior, donde la GNN produce $\theta_{\text{pred}}$ para $N$ no vistos), aquí se evalúa una estrategia más simple: transferir directamente los parámetros óptimos $\theta^*(N_1)$ como inicialización del VQE a $N_2 \neq N_1$, sin emplear el predictor MPNN.

\textbf{Tipo de experimento:} Simulación noiseless MPS ($\chi = 64$, $N = 40$--$80$, $p = 1$, cadena 1D, 3 topologías, COBYLA). Se evalúa si inicializar el VQE a $N = 50$ con $\theta^*(N = 40)$ mejora la convergencia frente a la inicialización estándar.

\textbf{Resultado:} Con solo 2 parámetros ($p = 1$), COBYLA \citep{singh2025} converge al mínimo global independientemente de la inicialización en 19--38 iteraciones. La transferencia no ofrece mejora medible porque el paisaje energético del HVA $p = 1$ para TFIM en la fase paramagnética es convexo con un único mínimo ---trivial para cualquier optimizador sin gradiente. El resultado se reproduce en las 3 topologías evaluadas (chain\_1d, heavy\_hex, square).

\textbf{Implicación:} La transferencia directa de parámetros VQE entre tamaños solo aporta valor cuando el número de parámetros es suficientemente grande para que el paisaje tenga mínimos locales ($\geq 10$ parámetros, régimen bond-resolved descrito en §\ref{subsec:bond_resolved}). Para $p = 1$--$2$ (2--4 parámetros), la inicialización es irrelevante ---la MPNN aporta valor únicamente como predictor zero-shot, no como proveedor de warm-start.

\section{Extensibilidad a Hamiltonianos de Ising}

Se valida la capacidad del framework para operar con variantes del Hamiltoniano de Ising más allá del TFIM estándar.

\subsection{TFIM con Campo Longitudinal}

El Hamiltoniano $H_{\text{TFIM+long}}$ (Ec.~\ref{eq:tfim_long}) añade un campo longitudinal $g \cdot Z_i$ que rompe la simetría $\mathbb{Z}_2$ explícitamente. El circuito HVA extendido incorpora rotaciones $R_Z(\theta_z)$ sin compuertas de dos qubits adicionales.

\begin{table}[h]
\centering
\caption{TFIM + campo longitudinal: fidelidad media del HVA extendido ($N = 6$, $p = 2$, 3 semillas, 3 topologías). Ref: \citet{dutta2015}.}
\label{tab:tfim_long}
\begin{tabular}{ccccc}
\toprule
\textbf{$g$} & \textbf{HVA estándar} & \textbf{HVA extendido} & \textbf{Mejora} & \textbf{CZ gates} \\
\midrule
0.0 & 0.990 & 0.991 & +0.001 & 10 \\
0.1 & 0.889 & 0.984 & +0.095 & 10 \\
0.2 & 0.778 & 0.985 & +0.207 & 10 \\
0.3 & 0.688 & 0.985 & +0.297 & 10 \\
0.5 & 0.556 & 0.987 & +0.431 & 10 \\
\bottomrule
\end{tabular}
\end{table}

El HVA extendido mantiene fidelidad $\geq 0.98$ hasta $g = 0.5$ con \textbf{cero compuertas de dos qubits adicionales}, demostrando que la extensión a campos de Ising más generales es viable sin costo en circuito.

Adicionalmente, el pipeline completo con TFIM longitudinal ($g = 0.3$) se validó para $N = 10$ sobre las 5 topologías con $p = 2$--$4$ (21 runs). Los resultados muestran un rendimiento comparable al TFIM estándar:

\begin{table}[h]
\centering
\caption{TFIM longitudinal: mejores configuraciones por topología ($N = 10$, $g = 0.3$, 40 h-points, pipeline completo Fases 1--3).}
\label{tab:tfim_long_deploy}
\begin{tabular}{lcccc}
\toprule
\textbf{Topología} & \textbf{Mejor $p$} & \textbf{PassRate} & \textbf{$\bar{F}_{\text{VQE}}$} & \textbf{Speedup} \\
\midrule
Heavy-Hex & 3 & \textbf{90\%} (35/39) & 0.995 & 82$\times$ \\
Cadena 1D & 3 & 92\% (36/39) & 0.997 & 62$\times$ \\
Escalera & 4 & 79\% (31/39) & 0.937 & 163$\times$ \\
Cuadrada & 4 & 67\% (26/39) & 0.927 & 191$\times$ \\
Triangular & 3 & 46\% (18/39) & 0.805 & 91$\times$ \\
\bottomrule
\end{tabular}
\end{table}

El hallazgo clave es que el TFIM longitudinal alcanza rendimiento equivalente al TFIM estándar (90\% vs 90\% en heavy-hex $p = 3$), confirmando que la extensión del Hamiltoniano no degrada el pipeline cuando los nuevos términos se implementan con compuertas de un qubit.

La Figura~\ref{fig:heatmap_passrate} visualiza la interacción entre tamaño de sistema ($N$) y profundidad del circuito ($p$) para el TFIM longitudinal en topología heavy-hex. Se observa que $p = 3$ es el sweet spot: alcanza 100\% de pass rate a $N = 16$ y $N = 20$, mientras que $p = 2$ funciona bien a $N \leq 10$ (85--88\%) pero se degrada a tamaños mayores. El resultado confirma que la profundidad moderada ($p = 3$) es suficiente para la extensión longitudinal incluso a escalas donde $p = 1$--$2$ es insuficiente.

\begin{figure}[h]
\centering
\includegraphics[width=0.75\textwidth]{thesis_plots/heatmap_pass_rate.png}
\caption{Mapa de calor del deploy pass rate ($\Delta E/\text{gap} < 5\%$) para TFIM longitudinal en topología heavy-hex. Ejes: tamaño de sistema $N$ vs profundidad del circuito $p$. Verde = pipeline funciona ($\geq 85\%$), amarillo = marginal (50--60\%), blanco = no evaluado. El régimen óptimo es $p = 3$, que alcanza 100\% para $N = 16$--$20$.}
\label{fig:heatmap_passrate}
\end{figure}

\subsection{TFIM Frustrado ($J_1$-$J_2$)}

El modelo frustrado alcanza fidelidad $\geq 0.99$ en simulación a $N = 6$, $p = 2$, $J_2 = 0.5$ (experimento T1c, 3 semillas, 100\% de acuerdo con la detección de fase por gradientes de pesos \citep{hernandes2025}), pero requiere 27 CZ gates a $N = 6$. \citet{teoh2025} validaron recientemente la simulación de un TFIM frustrado 2D con interacciones a segundos vecinos en un computador cuántico de iones atrapados, confirmando la viabilidad del enfoque para hardware de iones, aunque el costo en compuertas CZ excede la viabilidad práctica para procesadores superconductores.

\subsection{Heisenberg: Límite Fundamental de Expresividad}

Se evalúa exhaustivamente el modelo Heisenberg como caso negativo para documentar la frontera del HVA. Se ejecutaron \textbf{46 runs} independientes cubriendo dos variantes (Heisenberg estándar XXZ y Heisenberg transversal), 5 topologías, y profundidades $p = 1$--$4$ para $N = 10$. El resultado es concluyente: ninguna configuración alcanza PassRate $> 10\%$ y la fidelidad VQE permanece por debajo de la utilidad ($< 0.56$) en todos los casos.



Hallazgos principales de la evaluación exhaustiva:

\begin{enumerate}
    \item El fallo es \textbf{independiente} de profundidad ($p = 1$--$4$), topología (5 evaluadas), semilla y número de reinicios, confirmando un límite de expresividad del ansatz ---no de optimización.
    \item La variante transversal ($H = J(XX + YY + 0.5 \cdot ZZ) - h \cdot X$) muestra mejores fidelidades (0.54 vs 0.11) pero sigue siendo insuficiente para generar datos de entrenamiento válidos para la MPNN (umbral: fidelidad $\geq 0.93$).
    \item Incrementar de $p = 3$ a $p = 4$ mejora la fidelidad marginalmente (+0.14 en heavy-hex) pero duplica el costo computacional sin beneficio funcional.
    \item La MPNN alcanza MSE bajas en algunos runs Heisenberg (MSE $< 5 \times 10^{-4}$) porque el paisaje VQE \emph{es suave} ---simplemente converge a estados incorrectos. Esto demuestra que MSE bajo $\neq$ buena predicción cuando los datos de entrenamiento son erróneos.
\end{enumerate}

Este resultado es consistente con \citet{wiersema2020}, quienes demuestran que HVA para Heisenberg requiere $p \propto N$ capas, y con \citet{huang2026frustration}, quienes formalizan las limitaciones de expresividad inducidas por frustración geométrica en algoritmos variacionales.

\textbf{Conclusión sobre extensibilidad:} El framework GNN-HVA es directamente aplicable al TFIM estándar y sus extensiones de campo (longitudinal, frustrado), pero NO al modelo de Heisenberg con profundidades moderadas ($p \leq 4$). Esta es una limitación del ansatz, no del pipeline.

\subsection{Cadena de Kitaev: Incompatibilidad con el Framework}

\textbf{Tipo de experimento:} Simulación noiseless (\texttt{StatevectorEstimator}, $N = 4$--$6$, $p = 1$, 15 reinicios). Se evalúa la viabilidad del modelo de Kitaev ---cadena superconductora con pairing $p$-wave--- como candidato para demostrar detección de fases topológicas.

El Hamiltoniano de la cadena de Kitaev, tras transformación de Jordan-Wigner, contiene términos $XX + YY$ (hopping) y $XY - YX$ (pairing) que requieren compuertas $R_{XX}$ y $R_{YY}$ por bond. La evaluación revela tres barreras simultáneas que impiden su implementación en el framework:

(1) el presupuesto CX se duplica respecto al TFIM (20 CZ a $N=6$ vs. 10);
(2) el estado inicial $|+\rangle^{\otimes N}$ tiene overlap exponencialmente pequeño con el ground state de Kitaev (estado BCS de fermiones pareados);
(3) la fidelidad máxima alcanzada es 16\% incluso con 15 reinicios a $N = 4$ ---insuficiente para generar datos de entrenamiento para la MPNN (umbral: 93\%).

Estos resultados complementan el hallazgo para Heisenberg: modelos con interacciones de dos cuerpos adicionales (XX+YY) exceden la expresividad del HVA para profundidades moderadas.

\subsection{Resumen de Viabilidad por Modelo}

La Tabla~\ref{tab:candidates} sintetiza los resultados de extensibilidad. Los criterios de viabilidad son: expresividad del HVA (fidelidad $\geq 0.93$), observables locales disponibles, física no trivial a $N = 6$--$10$, presupuesto de compuertas CX compatible con mitigación, y valor científico distinto al TFIM estándar.

\begin{table}[h]
\centering
\caption{Evaluación de Hamiltonianos candidatos para el framework GNN-HVA. Solo TFIM + longitudinal pasa todos los criterios. Heisenberg y Kitaev se verificaron empíricamente}
\label{tab:candidates}
\begin{tabular}{lcccccl}
\toprule
\textbf{Modelo} & \textbf{Fid} & \textbf{Obs.} & \textbf{$N$} & \textbf{CX} & \textbf{Valor} & \textbf{Veredicto} \\
\midrule
TFIM + longitudinal & \checkmark & \checkmark & \checkmark & \checkmark & \checkmark & \textbf{Viable} \\
TFIM frustrado $J_1$-$J_2$ & $\sim$ & \checkmark & $\sim$ & $\times$ & \checkmark & No viable (CX) \\
Heisenberg XXZ & $\times$ & \checkmark & \checkmark & $\times$ & \checkmark & No viable (expresividad) \\
Kitaev chain & $\times$ & \checkmark & \checkmark & $\times$ & \checkmark & No viable (3 barreras) \\
\bottomrule
\end{tabular}
\end{table}

El resultado delimita la extensibilidad del framework: funciona para Hamiltonianos cuyos términos nuevos se mapean a compuertas de un qubit (rotaciones $R_X$, $R_Z$), sin incrementar el presupuesto de compuertas de dos qubits \citep{wiersema2020}. Modelos que requieren interacciones de dos cuerpos adicionales (XX, YY) quedan fuera del régimen viable con profundidades moderadas ($p \leq 4$).

\section{Aproximaciones Alternativas y Resultados Negativos}

Esta sección documenta técnicas y enfoques alternativos evaluados durante el desarrollo del pipeline que no produjeron mejoras o revelaron limitaciones fundamentales. Estos resultados negativos delimitan el espacio de aplicabilidad del framework y previenen la repetición de enfoques infructuosos.

\subsection{Extracción de Exponente Crítico $\nu$ desde Pesos de la Red (S8/S8b)}

\textbf{Tipo de experimento:} Simulación noiseless (StatevectorEstimator, $N = 4$, $6$, $8$, $10$; cadena 1D, $p = 2$; 5 semillas $\times$ 4 tamaños $= 20$ configuraciones por variante). Se evalúa si la posición del pico de gradiente de pesos $\|\partial W / \partial h\|$ sigue la ley de escalamiento finito $h_{\text{peak}}(N) = h_c + a \cdot N^{-1/\nu}$ con $\nu = 1$ (clase de universalidad del TFIM).

Se probaron dos variantes: S8 (MLP, entrada escalar $h$) y S8b (MPNN GINConv, entrada grafo de $N$ nodos). Ambas fallan:

\begin{table}[h]
\centering
\label{tab:s8_scaling}
\begin{tabular}{lcccl}
\toprule
\textbf{Variante} & \textbf{$h_{\text{peak}}$ mediana} & \textbf{$\nu_{\text{fit}}$} & \textbf{Dependencia en $N$} \\
\midrule
S8 (MLP) & 0.704 (fijo $\forall N$) & 5.0 & Ninguna \\
S8b (MPNN) & 0.500 (borde del rango) & 5.0 & Ninguna \\
\bottomrule
\end{tabular}
\caption{Escalamiento finito del pico de gradiente de pesos: S8 (MLP) y S8b (MPNN). $\nu_{\text{fit}}$ debería ser $\approx 1.0$ para TFIM; se obtiene 5.0 (límite superior del ajuste).}

\end{table}

\textbf{Causa raíz:} La señal de gradiente está dominada por la transición de calidad de los datos VQE en $h < 1.0$ (fuera del régimen válido del HVA $p = 2$). Esta transición de calidad de datos ---no la transición de fase física--- genera el pico más pronunciado del gradiente. El contraste con \citet{hernandes2025}, quienes extraen exitosamente fronteras de fase desde el espacio de pesos, se explica por tres diferencias metodológicas: (1) uso de estados fundamentales exactos (NQS) vs. nuestros parámetros VQE-optimizados, (2) métrica de distancia en espacio de pesos vs. norma de gradiente, y (3) calidad de datos uniformemente alta vs. degradación en la región crítica.

\textbf{Implicación:} El método D1 de detección de fases por gradientes de pesos es \emph{cualitativo} (detecta que una transición existe en $h \approx 0.7$--$1.2$) pero \emph{no cuantitativo} (no puede extraer $\nu$). Para extracción cuantitativa de exponentes críticos se requiere: (a) datos de ground truth exacto como entrada, (b) una métrica diferente (distancia en espacio de pesos, información de Fisher), o (c) restricción del análisis al régimen válido ($h > h_{\min}$).

\subsection{Dynamic Parameter Prediction No Mejora el Warm-Start (F1)}

\textbf{Tipo de experimento:} Simulación noiseless (StatevectorEstimator, $N = 6$, $p = 2$, cadena 1D, 3 semillas). Se evalúa si una técnica de predicción dinámica de parámetros (DyPP) ---que ajusta la inicialización punto a punto usando la información del gradiente previo--- reduce significativamente las iteraciones VQE más allá del warm-start descendente estándar \citep{puig2025}.

\textbf{Resultado:} DyPP produce un ahorro de solo 8--13\% en iteraciones (de $\sim$14 iteraciones con warm-start estándar a $\sim$12 con DyPP), con una tasa de aprobación del 64\%.

\textbf{Implicación:} El warm-start descendente ya proporciona inicialización casi-óptima para el HVA con 4 parámetros ---el paisaje energético es suficientemente suave \citep{mele2022} que no queda margen significativo de mejora mediante predicción adaptativa. Este resultado confirma que la simplificación del pipeline (warm-start puro, sin DyPP) es la estrategia correcta.

\section{Pruebas Sistemáticas de Funcionamiento}

\textbf{Tipo de experimento:} Meta-análisis sobre 174 ejecuciones completas del pipeline (simulación noiseless, múltiples topologías y tamaños). Se aplica diagnóstico automatizado post-hoc para clasificar las causas de fallo y cuantificar la efectividad del sistema de early-stopping.

Se presentan los resultados del diagnóstico automatizado, proporcionando evidencia sistemática del funcionamiento correcto.

\subsection{Análisis de Causa Raíz}

El diagnóstico automatizado sobre 174 ejecuciones clasificó 76 fallos por causa raíz. La tasa de resultados útiles (confirmados + negativos válidos) es del 84\% (Figura~\ref{fig:experiment_verdicts}).

\begin{figure}[h]
\centering
\includegraphics[width=0.90\textwidth]{tesis-figures/fig_failure_mode_treemap.pdf}
\caption{Análisis de causa raíz sobre 174 ejecuciones del pipeline. El 69\% de los fallos es prevenible: CHAIN\_BREAK ($\theta_{\text{smooth}} > 1.0$, detectable en Fase 2) y MPNN\_OVERFIT (gen\_gap $> 0.01$, detectable en Fase 3). El 23\% restante (BOUNDARY\_EFFECT + OUTSIDE\_REGIME) es prevenible en pre-ejecución con la regla de $h_{\min}$.}
\label{fig:failure_treemap}
\end{figure}

La relación entre calidad de la MPNN y resultado final del pipeline se visualiza en la Figura~\ref{fig:mpnn_scatter}: existe correlación log-log ($r = 0.46$) entre el generalization gap de la MPNN y el $\Delta E/\text{gap}$ final, confirmando que el gen\_gap $> 0.01$ es un predictor confiable de fallo.

\begin{figure}[t]
\centering
\includegraphics[width=0.85\textwidth]{tesis-figures/fig_mpnn_generalization_scatter.pdf}
\caption{Correlación entre la calidad de la MPNN (generalization gap) y la precisión final del pipeline ($\Delta E/\text{gap}$). El umbral gen\_gap $= 0.01$ separa eficazmente ejecuciones exitosas ($< 5\%$) de fallidas, justificando su uso como criterio de aborto temprano en Fase 3.}
\label{fig:mpnn_scatter}
\end{figure}

\subsection{Eficiencia de Datos}
\label{sec:eficiencia_datos}

La MPNN demuestra alta eficiencia de datos: logra $\Delta E/\text{gap} < 5\%$ con tan solo 5 puntos de entrenamiento para semillas favorables, y la recomendación conservadora es $k = 7$--$9$ puntos para robustez cross-seed. A $N = 10$ ($p = 3$, cadena 1D, 3 semillas), el $\Delta E/\text{gap}$ medio varía solo entre 2.94\% ($k = 17$, completo) y 3.37\% ($k = 5$, mínimo viable), confirmando que la función $h \mapsto \theta^*$ es suficientemente suave para aprenderse con pocas muestras. Para $N = 30$--$60$ ($p = 1$), 16--39 puntos de entrenamiento logran 100\% PassRate.

\section{Comparativa Sistemática con la Literatura}

Se realiza una comparación cuantitativa del pipeline GNN-HVA con los procedimientos más relevantes de la literatura reciente.

\textbf{Diferencias del enfoque propuesto respecto a los trabajos revisados:}

\begin{enumerate}
    \item \textbf{Escala de validación:} El pipeline se valida hasta $N = 20$ (simulación exacta) y $N = 40$ ($p = 1$, MPS, UnifiedMPNN), mientras que los trabajos revisados operan a $N \leq 12$ \citep{miao2024, zhang2025, zou2026}.

    \item \textbf{Eliminación de la optimización iterativa:} Los trabajos previos reducen iteraciones (Qracle: $-$64\% \citep{zhang2025}; Flow-VQE: $50\times$ \citep{zou2026}) pero mantienen un bucle de refinamiento. Nuestro pipeline logra predicciones sub-5\% sin refinamiento, gracias a la calidad de los datos de warm-start \citep{puig2025, schiffer2026}.

    \item \textbf{Validación multi-topología:} 5 topologías con hiperparámetros unificados --- aspecto no evaluado en los trabajos revisados.

    \item \textbf{Extensibilidad documentada:} TFIM + longitudinal (tasa de aprobación 90\%+), Heisenberg (46 runs, 0\% --- límite formal documentado, consistente con \citet{wiersema2020, huang2026frustration}).
\end{enumerate}

\begin{table}[h]
\centering
\caption{Comparativa del pipeline propuesto con métodos de la literatura para inicialización de parámetros VQE en sistemas de espines.}
\label{tab:comparison_literature}
\begin{tabular}{p{2.5cm}p{2cm}p{2cm}p{2cm}p{2.5cm}p{2cm}}
\toprule
\textbf{Método} & \textbf{Arquitectura} & \textbf{Sistema} & \textbf{$N$ máx} & \textbf{Reducción iter.} & \textbf{Multi-topo} \\
\midrule
NN-VQE \citep{miao2024} & MLP & Heisenberg & 10 & $\sim$5$\times$ & No \\
Qracle \citep{zhang2025} & GNN & Varios & 12 & $\sim$2.7$\times$ & Sí \\
Flow-VQE \citep{zou2026} & NF & H$_2$, LiH & 12 & 50$\times$ & No \\
Lee et al. \citep{lee2026} & GAE+NN & Ising & 8 & $\sim$3$\times$ & No \\
\midrule
\textbf{GNN-HVA} & \textbf{GINConv MPNN} & \textbf{TFIM (+ext)} & \textbf{200} & \textbf{29--500$\times$} & \textbf{Sí} \\
\bottomrule
\end{tabular}
\end{table}

\subsection{Comparación Cuantitativa: Warm-Start VQE vs. Predicción Directa}

La estrategia de warm-start descendente ---propuesta teóricamente por \citet{mele2022} y formalizada por \citet{puig2025}--- reduce las iteraciones de 500--1000 (inicialización aleatoria \citep{tilly2022}) a 20--50 por punto. La predicción MPNN, siguiendo el paradigma propuesto por \citet{miao2024} y \citet{zhang2025}, produce $\Delta E/\text{gap} < 5\%$ sin refinamiento en el 95\% de los casos (Tabla~\ref{tab:results_baseline} y Tabla~\ref{tab:results_p1}). La diferencia cuantitativa respecto a Qracle \citep{zhang2025} ---que reduce un 64\% las iteraciones pero mantiene refinamiento--- se atribuye a la mayor calidad de los datos de entrenamiento producidos por el warm-start descendente \citep{puig2025, schiffer2026}, que genera $\theta^*$ de alta fidelidad para entrenar la MPNN.

La Figura~\ref{fig:warmstart_vs_random} cuantifica esta ventaja en función del tamaño del sistema. Mientras el error del VQE con inicialización aleatoria crece exponencialmente con $N$ (superando el umbral del 5\% desde $N = 10$), la predicción MPNN se mantiene estable por debajo del 3\% para todos los tamaños evaluados. El factor de mejora crece de $2\times$ ($N = 6$) a $500\times$ ($N = 20$), demostrando que la ventaja del pipeline se amplifica con la escala del sistema ---precisamente el régimen donde la alternativa convencional se vuelve impracticable.

\begin{figure}[h]
\centering
\includegraphics[width=0.95\textwidth]{thesis_plots/warmstart_vs_random_v2.png}
\caption{Comparación cuantitativa entre predicción MPNN (zero-shot) y VQE con inicialización aleatoria (TFIM longitudinal, heavy-hex, $p = 3$). (a) $\Delta E/\text{gap}$ en escala logarítmica: la MPNN se mantiene bajo 5\% mientras Random-VQE diverge con $N$. (b) Factor de mejora en escala logarítmica: crece de $2\times$ a $500\times$. (c) Error absoluto $|\Delta E|$ decrece con $N$ mientras la fidelidad se mantiene $\geq 0.99$. (d) Tasa de victoria MPNN ($\geq 90\%$) y esfuerzo del VQE aleatorio (37--85 iteraciones por punto).}
\label{fig:warmstart_vs_random}
\end{figure}

\subsection{Ablación: ¿Es la MPNN Necesaria o Basta con Interpolación?}

Dado que los parámetros óptimos $\theta^*(h)$ varían suavemente dentro de una fase \citep{mele2022}, surge una pregunta legítima: ¿es necesario entrenar una red neuronal, o basta con interpolar linealmente entre los puntos de entrenamiento? Para responderla, se compara la MPNN contra un baseline trivial: interpolación lineal a tramos (\texttt{scipy.interpolate.interp1d}), que construye $\hat{\theta}(h) = \theta^*_i + (\theta^*_{i+1} - \theta^*_i) \cdot \frac{h - h_i}{h_{i+1} - h_i}$ entre cada par de puntos vecinos conocidos. Ambos métodos usan exactamente los mismos datos de entrenamiento; la diferencia es que la interpolación no requiere modelo ni entrenamiento.

\textbf{Conclusión.} Para el caso base (HVA con parámetros compartidos, 2--4 dimensiones, puntos de test dentro del rango de entrenamiento), la interpolación lineal es un baseline competitivo ---la MPNN no ofrece ventaja significativa en precisión. Esto es consistente con la suavidad del paisaje reportada por \citet{mele2022}: en baja dimensionalidad, $\theta^*(h)$ es esencialmente lineal dentro de una fase.

La MPNN demuestra su valor en dos escenarios donde la interpolación fracasa:
\begin{enumerate}
    \item \textbf{Alta dimensionalidad} (bond-resolved, 79 parámetros): interpolar en $\mathbb{R}^{79}$ con 45 puntos es geométricamente inviable; la MPNN generaliza con factor $4414\times$ sobre random.
    \item \textbf{Eficiencia de datos para $N$ grande}: a $N = 30$--$60$, la MPNN logra 100\% PassRate con 16--39 puntos de entrenamiento, donde la interpolación lineal requeriría muestreo denso en cada dimensión.
\end{enumerate}

\textbf{Limitación identificada:} la extrapolación cross-N (entrenar en $N_1$, predecir en $N_2 \neq N_1$) \emph{no funciona} para HVA con parámetros compartidos ($\Delta E/\text{gap} = 688\%$ media, $N = 10, 20, 40 \to N = 60$). La causa es que $\theta^*(h)$ depende de $N$ a través de la relación entre $h$ y la frontera de expresividad, y el \texttt{global\_mean\_pool} produce embeddings incompatibles entre tamaños distintos.

\section{Análisis de Límites Físicos}

\subsection{Frontera de Expresividad del HVA}

La frontera donde el HVA no puede representar el estado fundamental ---independientemente de la estrategia de optimización--- se localiza en $h \approx 1.25$ para $N = 6$ ($p = 2$) y $h \approx 1.3$ para $N = 10$ ($p = 3$). Esta frontera coincide con la región donde la entropía de entrelazamiento excede la capacidad del circuito. El resultado ha sido confirmado independientemente por \citet{tripathi2026} y es consistente con la cota teórica de \citet{sumeet2025} ($N/2$ capas para el límite termodinámico). Con profundidad suficiente ($p = 5$, $N = 8$), la frontera se desplaza hasta $h \approx 0.97 \approx h_c$ (véase Tabla~\ref{tab:scaling}), confirmando que el límite es de profundidad, no algorítmico.

\textbf{Comportamiento con $N$ (datos MPS determinista, $\chi = 64$, $N = 20$--$250$):} La frontera $h_{\min}(N, p)$ exhibe dos regímenes cualitativamente distintos según la profundidad del circuito:

\begin{itemize}
    \item Para $p = 1$--$2$: $h_{\min}$ crece linealmente con $N$, con pendientes pequeñas. Los ajustes empíricos son $h_{\min}(p{=}1) = 2.36 + 0.0073 \cdot N$ ($R^2 = 0.91$, 10 puntos, $N = 20$--$250$) y $h_{\min}(p{=}2) = 1.57 + 0.0050 \cdot N$ ($R^2 = 0.95$, 8 puntos, $N = 20$--$120$).
    \item Para $p \geq 3$: $h_{\min}$ es \emph{cuasi-constante}, con $h_{\min}(p{=}3) \approx 1.6 \pm 0.1$ ($R^2 = 0.46$ para un modelo lineal --- inapropiado) y $h_{\min}(p{=}4) \approx 1.4 \pm 0.1$ ($R^2 = 0.73$, pendiente no significativa).
\end{itemize}

Este hallazgo es consistente con la \emph{area law} de entrelazamiento en 1D: a $h/h_c$ fijo, la entropía de entrelazamiento es independiente de $N$ para estados fundamentales gapped. Si el circuito HVA posee suficiente profundidad ($p \geq 3$) para capturar dicho entrelazamiento, la frontera no se desplaza al escalar $N$. La varianza residual entre semillas ($\sigma \approx 0.10$--$0.13$) domina sobre cualquier tendencia con $N$ para $p \geq 3$.

\textbf{Dependencia con el modelo:} Para el TFIM, la frontera decrece con $p$ y converge a $h_c$ ---pero esto es específico de la coincidencia entre la estructura del HVA ($R_{ZZ} + R_X$) y los términos del Hamiltoniano ($ZZ + X$). Para el modelo Heisenberg XXZ ($XX + YY + \Delta \cdot ZZ$), la frontera se estabiliza en $h_{\min} \approx 3.5$ \emph{independientemente} de $p$ hasta $p = 8$ ($N = 10$, rango 3.67--3.50): el HVA con estructura TFIM no puede capturar las correlaciones XX/YY requeridas. Este resultado confirma que la expresividad del ansatz es específica al Hamiltoniano \citep{wiersema2020}, y que la convergencia $p \to N/2$ de \citet{sumeet2025} aplica solo cuando circuito y Hamiltoniano comparten estructura de términos.

\textbf{Dependencia con la topología:} El número de coordinación máximo ($z_{\max}$) es el factor dominante. A $N = 10$ y profundidad óptima: cadena 1D ($z = 2$): $h_{\min} = 1.09$; heavy-hex ($z = 3$): $h_{\min} = 1.12$ --- funcionalmente equivalente a la cadena, excelente para hardware IBM; escalera ($z = 3$): $h_{\min} = 1.67$; cuadrada ($z = 4$): $h_{\min} = 1.68$; triangular ($z = 6$): $h_{\min} = 2.20$, con una pared estructural en $\Delta E/\text{gap} \approx 0.5$ por debajo de $h = 2$ que constituye un límite del ansatz con parámetros globales, no un problema de optimización. Ninguna topología cruza $h_c = 1.0$ con parametrización uniforme --- únicamente la parametrización bond-resolved lo logra ($h_{\min} = 0.83$, §\ref{subsec:bond_resolved}).

\textbf{Limitación computacional a $p \geq 4$:} Para profundidades $p \geq 4$ con $N > 40$ y $h$ cerca de $h_c$, la obtención del ground truth vía DMRG requiere dimensión de enlace $\chi > 200$, excediendo los recursos prácticos del simulador MPS. Esto impide medir $h_{\min}$ para $p = 4$ a $N = 100$--$120$. Sin embargo, los datos disponibles ($N = 20$--$80$) ya confirman que $h_{\min}(p{=}4) \approx 1.4$ es cuasi-constante, y el crecimiento leve observado (1.36 a $N = 20$ $\to$ 1.51 a $N = 80$) se atribuye a la dificultad creciente del paisaje de optimización VQE --- no a un artefacto numérico del MPS, ya que los resultados son exactos a machine epsilon.

\subsection{Detección de Fases en el Espacio de Parámetros}

El pipeline produce como subproducto el conjunto de parámetros óptimos $\theta^*(h)$ para cada punto del barrido. Si la transición de fase cuántica induce un cambio estructural en estos parámetros, la transición puede detectarse sin mediciones cuánticas adicionales. Se evalúan dos indicadores complementarios.

Siguiendo \citet{hernandes2025}, se analizan los gradientes de los pesos de la MPNN respecto a $h$. Se detectan picos en $\|\partial W / \partial h\|$ en la región $h \approx 1.0$--$1.2$, indicando que la red codifica la transición de fase internamente.

Complementariamente, se aplica Análisis de Componentes Principales (PCA) al dataset de parámetros óptimos $\{\theta^*(h_i)\}_{i=1}^{n}$, tratando cada vector $\theta^* \in \mathbb{R}^{2p}$ como una observación y cada valor de $h$ como una muestra. PCA identifica las direcciones de máxima varianza en este espacio: si una única componente principal (PC1) captura la mayor parte de la varianza total, entonces la trayectoria $\theta^*(h)$ se confina a una curva esencialmente unidimensional, y las discontinuidades de esta curva señalan transiciones de fase.

La Tabla~\ref{tab:unsupervised_detection} compara tres métodos de detección no supervisada:

\begin{table}[h]
\centering
\caption{Detección no supervisada de la transición de fase ($N = 6$, cadena 1D, $p = 2$).}
\label{tab:unsupervised_detection}
\begin{tabular}{lccc}
\toprule
\textbf{Método} & \textbf{Pico detectado ($h$)} & \textbf{$\Delta h$ vs $h_c = 1.0$} \\
\midrule
Gradiente de pesos MPNN & 0.99 & 0.01 \\
PCA de $\theta_{\text{opt}}(h)$ & 1.25 & 0.25 \\
$|\partial\theta/\partial h|$ & 1.25 & 0.25 \\
\bottomrule
\end{tabular}
\end{table}

El resultado del PCA confirma que el espacio de parámetros HVA es efectivamente unidimensional: PC1 explica el 99.96\% de la varianza total, lo que significa que la variación de $\theta^*(h)$ al recorrer el diagrama de fases se confina a una única dirección en $\mathbb{R}^{2p}$. A $N = 100$ con un grid de $h$ que cubre la región crítica, el pico de varianza de la proyección sobre PC1 converge a $h = 1.033 \pm 0.047$, consistente con $h_c = 1.0$ \citep{dutta2015}. Sin embargo, el método es cualitativo: no puede extraer el exponente crítico $\nu$ (véase resultados negativos S8/S8b).

Este resultado confirma que la detección PCA es un método válido de costo cero cuántico, pero requiere datos VQE que abarquen la región crítica ---condición satisfecha solo para $N \leq 10$ bajo la restricción de régimen válido ($h_{\min}$). Cabe señalar que \citet{yin2025} han propuesto independientemente un enfoque más sofisticado para la misma tarea: un autoencoder variacional (VAE) con mecanismo de atención entrenado sobre parámetros de circuitos variacionales para detectar QPTs de manera no supervisada. El método de Yin et al.\ es más poderoso que nuestro PCA ---captura relaciones no lineales en el espacio de parámetros--- pero también más complejo. Nuestro análisis PCA, siendo lineal y de costo computacional negligible, ofrece un diagnóstico de primer orden suficiente para confirmar cualitativamente la presencia de la transición.

\section{Discusión: ¿Por qué Funciona la Metodología?}

La eficacia del pipeline GNN-HVA se fundamenta en la convergencia de cinco propiedades teóricas y empíricas que simplifican radicalmente el problema de aprendizaje:

\textbf{1. El paisaje del HVA es suave dentro de una fase.} \citet{mele2022} demostraron que $\theta^*(h)$ varía continuamente con $h$ para Hamiltonianos parametrizados con gap espectral. Para TFIM con HVA global ($p = 1$), la relación $h \mapsto \theta^*$ es una función de solo 2 parámetros ($\theta_{ZZ}$, $\theta_X$) que varía monótonamente con $h$ en el régimen paramagnético. Cuantitativamente, la métrica $\theta_{\text{smooth}} = \max_i \|\theta^*(h_i) - \theta^*(h_{i-1})\|_\infty$ da valores de 0.01--0.05 en el régimen válido ---sin saltos, discontinuidades ni oscilaciones rápidas. Esto implica que la función $h \mapsto \theta^*$ es altamente predecible: pocas decenas de puntos de entrenamiento bastan para interpolar con precisión suficiente. Nuestros resultados de eficiencia de datos (5 puntos bastan para semillas favorables, Sección~\ref{sec:eficiencia_datos}) confirman empíricamente esta suavidad.

\textbf{2. La física restringe el espacio de soluciones y la dimensionalidad del output es mínima.} El Hamiltoniano TFIM tiene solo dos fases (paramagnética y ferromagnética) con un crossover suave. En el régimen $h \gg h_c$ ---donde opera la mayoría de las configuraciones validadas--- el estado fundamental está próximo a $|+\rangle^{\otimes N}$, y el HVA solo necesita aplicar una corrección perturbativa. Los parámetros óptimos están concentrados en una región compacta ($|\theta^*| < 0.5$ típicamente). Además, con HVA de parametrización global y $p = 1$, el output de la MPNN son solo 2 números ($\theta_{ZZ}$, $\theta_X$). Una red con $\sim$99K parámetros (3 capas GINConv, hidden\_dim$=128$) que mapea un grafo de $N$ nodos a 2 valores escalares dispone de enorme capacidad relativa ---incluso con 9 puntos de entrenamiento, el modelo tiene suficiente señal para capturar la relación funcional, especialmente cuando los puntos están distribuidos a lo largo de una curva suave y los labels no tienen ruido (el VQE converge a $< 0.5\%$ de error).

\textbf{3. La estructura de grafo informa la predicción.} El mecanismo de paso de mensajes de GINConv captura la topología del Hamiltoniano ---qué qubits interactúan y con qué intensidad--- sin codificación manual de features. Esto explica por qué topologías con mayor conectividad (heavy-hex, escalera) producen mejores predicciones: más aristas $=$ más canales de información para la GNN. \citet{slavin2025} confirman independientemente que la geometría de la red es suficiente para predecir magnetización en sistemas de Ising.

\textbf{4. La inicialización informada elimina los mínimos locales y multiplica datos.} El warm-start descendente, combinado con la suavidad del paisaje HVA, garantiza que la optimización converge al mínimo global en cada punto del barrido. Esto produce datos de entrenamiento de alta calidad para la MPNN ---a diferencia de inicializaciones aleatorias que frecuentemente quedan atrapadas en mínimos locales. \citet{puig2025} formalizan este argumento: el warm-start proporciona varianzas de pérdida mayores (gradientes más fuertes), asegurando convergencia confiable. Adicionalmente, la ejecución multi-semilla multiplica los datos sin costo de exploración adicional: con 3 semillas, 9 valores de $h$ producen 27 puntos de entrenamiento, cada uno con una $\theta^*$ ligeramente distinta (ruido de optimización $\sim$0.4--0.8\% $\Delta E/\text{gap}$ entre semillas) pero igualmente válida. Esto proporciona a la MPNN una noción implícita de la distribución de soluciones, mejorando la robustez de la predicción. Cabe notar que este argumento tiene un límite con el tamaño del sistema: a $N = 20$ (experimento S3), el paisaje exhibe 2--3 mínimos locales distintos y condiciones de curvatura más difíciles ($\kappa$ variable entre semillas), requiriendo $\geq 7$ reinicios para explorar todas las cuencas de atracción ---a diferencia de $N \leq 10$ donde un único mínimo global domina el paisaje.

\textbf{5. La profundidad moderada es compatible con hardware futuro.} El resultado de \citet{mele2026} establece que la profundidad efectiva útil escala como $O(\log n)$ en presencia de ruido. Los circuitos HVA de $p = 1$--$4$ capas empleados en este trabajo satisfacen holgadamente esta cota para $N = 4$--$20$. En simulación ideal, estos circuitos son suficientes para resolver la física fuera de la región crítica ($h \geq h_{\min}$), como confirman los resultados de las Tablas~\ref{tab:cross_topo} y \ref{tab:cross_topo_depth}.

La combinación de estos cinco elementos produce un pipeline que funciona de manera confiable: la física suave del TFIM fuera de la región crítica + la baja dimensionalidad del output + la estructura de grafo informativa + la inicialización de calidad + el ansatz apropiado convergen en un sistema donde la MPNN solo necesita capturar una curva suave de pocos valores ---un problema de aprendizaje fundamentalmente distinto de un ML genérico donde la relación input$\to$output puede ser arbitrariamente compleja.

La Figura~\ref{fig:gap_widens} sintetiza visualmente la consecuencia práctica de estas propiedades: mientras la MPNN mantiene un error esencialmente constante al crecer $N$, el VQE con inicialización aleatoria diverge exponencialmente. La ventaja del pipeline no es marginal ---se amplifica con la escala del sistema, alcanzando un factor $500\times$ a $N = 20$.

\begin{figure}[h]
\centering
\includegraphics[width=0.85\textwidth]{thesis_plots/gap_widens_mpnn_vs_random.png}
\caption{Divergencia de la ventaja MPNN con el tamaño del sistema (TFIM longitudinal, heavy-hex, $p = 3$). La predicción MPNN (verde) se mantiene estable bajo el umbral de 5\% para todo $N$, mientras que el VQE con inicialización aleatoria (rojo) diverge exponencialmente. La zona sombreada azul representa la ventaja acumulada del pipeline, que crece de $2\times$ ($N = 6$) a $500\times$ ($N = 20$).}
\label{fig:gap_widens}
\end{figure}

\section{Resultados del Resumen de Experimentos}

El análisis de la campaña experimental muestra que la tasa de resultados útiles (confirmados + rechazados con resultado negativo válido) es del 84\% (41/49), demostrando la madurez de la metodología experimental. Los 8 experimentos rechazados producen conocimiento accionable: delimitan la aplicabilidad del framework (cross-topology transfer falla con $\Delta E/\text{gap} = 625\%$--$719\%$ para transferencia bidireccional triangular $\leftrightarrow$ heavy-hex, Heisenberg falla, escalamiento de exponente crítico falla, DyPP redundante).

\begin{figure}[h]
\centering
\includegraphics[width=0.80\textwidth]{tesis-figures/fig_experiment_verdicts_overview.pdf}
\caption{Distribución de veredictos de la campaña experimental por categoría. De 49 experimentos formales, 33 confirman su hipótesis, 8 producen resultados negativos válidos (contribuciones a los límites del framework), y 8 fallan genuinamente.}
\label{fig:experiment_verdicts}
\end{figure}

\section{Ventaja de la Solución y Aplicabilidad}

La evaluación sistemática demuestra que el pipeline GNN-HVA resuelve el problema planteado ---clasificación de fases cuánticas con costo cuántico mínimo--- con las siguientes ventajas cuantificables respecto a los enfoques revisados en la literatura:

\begin{enumerate}
    \item \textbf{Aceleración $29$--$500\times$:} La MPNN predice parámetros en una única pasada forward clásica, eliminando el bucle iterativo de VQE. Frente a métodos como Qracle ($-$64\% iteraciones \citep{zhang2025}), el pipeline logra predicciones sub-5\% sin refinamiento iterativo, gracias a warm-start de alta calidad \citep{puig2025} y GINConv \citep{xu2019}.

    \item \textbf{Multi-topología con hiperparámetros unificados:} Un único conjunto funciona en 5 topologías (Tabla~\ref{tab:cross_topo}), aspecto no evaluado en trabajos previos.

    \item \textbf{Escalabilidad:} El pipeline opera desde $N = 4$ hasta $N = 20$ con simulación exacta (Tabla~\ref{tab:scaling}), y se extiende a $N = 40$ con la variante UnifiedMPNN manteniendo $|\Delta E|/N \approx$ constante (Tabla~\ref{tab:large_n_chain}), confirmando que la predicción escala extensivamente con el tamaño del sistema dentro del régimen válido.
\end{enumerate}

\textbf{Aplicabilidad en materia condensada.} El pipeline se posiciona como herramienta metodológica para la simulación cuántica variacional de sistemas de espines:

\begin{itemize}
    \item \textbf{Screening de diagramas de fases:} Un experimentalista que disponga de un procesador cuántico puede entrenar la MPNN con 16--39 puntos VQE y predecir el diagrama completo sin optimización adicional, reduciendo el uso de QPU de días a minutos para sistemas con paisaje energético suave.

    \item \textbf{Clasificación sin reconstrucción de estado:} La filosofía de observables locales ($\langle X_i \rangle$, $\langle Z_i Z_{i+1} \rangle$) es compatible con hardware ruidoso, donde la tomografía completa es prohibitiva \citep{sharma2026}. El pipeline identifica la fase con mediciones de bajo costo.

    \item \textbf{Benchmark de pipelines híbridos:} Los criterios cuantitativos ($\Delta E/\text{gap} < 5\%$, fidelidad $\geq 0.93$, tasa de aprobación por configuración) proporcionan una metodología de evaluación reproducible para validar ansätze o predictores alternativos.

    \item \textbf{Generación eficiente de datos de entrenamiento:} El warm-start reduce el costo de generar datasets VQE en $10$--$50\times$ \citep{puig2025}, habilitando conjuntos de entrenamiento a gran escala para ML cuántico.
\end{itemize}

\textbf{Nota sobre simulabilidad clásica.} A las escalas validadas ($N \leq 20$, 1D), todos los resultados son reproducibles clásicamente ---la frontera de ventaja cuántica se sitúa en $N \approx 20$ para 2D \citep{martin2026}. El trabajo demuestra la validez del \emph{método} en un régimen verificable; la ventaja cuántica se materializa al escalar a sistemas 2D frustrados donde DMRG pierde eficiencia.

\textbf{Limitaciones:} Modelos con interacciones XX/YY (Heisenberg, Kitaev) hasta $p = 4$ (46 runs, 0\% PassRate, Sección~\ref{sec:extensibilidad}); la región crítica exacta ($h \approx h_c$) donde el HVA carece de expresividad \citep{tripathi2026, sumeet2025}; y sistemas con transiciones puramente topológicas.

\section{Código Desarrollado}
\label{sec:codigo}

El pipeline GNN-HVA se implementa como un paquete Python instalable (\texttt{qmbp\_simulation}) con arquitectura modular y control de calidad automatizado. El código fuente está disponible públicamente en:

\begin{center}
\url{https://github.com/franraineri/qmbp_gnn_adapt-vqe}
\end{center}

\textbf{Estructura del repositorio:}
\begin{itemize}
    \item \texttt{src/qmbp\_simulation/} --- Paquete principal (12 módulos: models, solvers, circuits, execution, optimizers, predictors, pipeline, analysis, framework, utils).
    \item \texttt{tests/} --- tests unitarios y de integración (pytest).
    \item \texttt{scripts/experiment\_runners/} --- Scripts reproducibles para cada experimento.
    \item \texttt{results/} --- Resultados en formato JSON.
    \item \texttt{documentation/} --- Bitácoras, análisis y figuras.
\end{itemize}

\textbf{Características técnicas:}
\begin{itemize}
    \item Compatible con Qiskit 2.x (Primitives V2, \texttt{SparsePauliOp}).
    \item Backend configurable: statevector ($N \leq 22$), MPS ($N > 22$), hardware IBM.
    \item MPNN basada en PyTorch Geometric \citep{fey2019} (GINConv, 3 capas).
    \item Validación automática de resultados VQE (principio variacional, convergencia).
    \item Validación de predicciones $\theta$ (7 niveles configurables: bounds, NaN, interpolación, fidelidad, gradiente, sensibilidad).
    \item Diagnóstico automatizado de fallos con clasificación de causa raíz.
    \item Generación automatizada de figuras para la tesis (\texttt{make figures-thesis}).
\end{itemize}


%---------------------------
% CAPÍTULO 6: CONCLUSIONES Y TRABAJO FUTURO
%---------------------------

\chapter{Conclusiones y Trabajo Futuro}

\section{Resumen del Problema y Enfoque}

El presente trabajo abordó el problema de la caracterización de fases cuánticas, donde las tres barreras fundamentales ---barren plateaus \citep{mcclean2018}, truncamiento por ruido \citep{mele2026} y costo de evaluación \citep{tilly2022}--- limitan la aplicación directa de VQE convencional. Se desarrolló una Arquitectura Predictiva Híbrida (GNN-HVA) que elimina la optimización cuántica iterativa mediante la predicción clásica de parámetros variacionales ---siguiendo el paradigma propuesto por \citet{miao2024} y \citet{zhang2025}---, operando dentro del régimen donde el HVA es simultáneamente entrenable \citep{cerezo2021} y expresivo \citep{wiersema2020}. La contribución específica de este trabajo es la \emph{integración} de estas técnicas existentes en un pipeline end-to-end y su \emph{validación sistemática} mediante simulación ideal sobre múltiples topologías ($5$), profundidades ($p = 1$--$4$), modelos (TFIM + variantes) y escalas ($N = 4$--$20$, 400+ ejecuciones).

La solución se considera válida porque:
\begin{itemize}
    \item Satisface el criterio de éxito primario ($\Delta E/\text{gap} < 5\%$) en el 100\% de los puntos dentro del régimen válido para configuraciones optimizadas (Tablas~\ref{tab:cross_topo}--\ref{tab:scaling}).
    \item Se fundamenta en resultados teóricos rigurosos: transferibilidad de parámetros \citep{mele2022, puig2025}, expresividad máxima de GINConv \citep{xu2019}, y validación independiente del HVA para TFIM \citep{tripathi2026}.
    \item Los límites observados son físicos (expresividad del ansatz), no algorítmicos, confirmados independientemente en la literatura.
\end{itemize}

\section{Contribuciones y Cumplimiento de Objetivos}

Se resumen las contribuciones del trabajo y su relación con los objetivos planteados:

\begin{enumerate}
    \item \textbf{El pipeline GNN-HVA clasifica correctamente fases cuánticas} con $\Delta E/\text{gap} < 5\%$ en el régimen operativo válido, validado sobre 400+ ejecuciones en simulación ideal, 5 topologías, profundidades $p = 1$--$4$, y tamaños de sistema de $N = 4$ a $N = 20$ (Tablas~\ref{tab:cross_topo}--\ref{tab:cross_topo_depth}). La tasa de aprobación alcanza 95--100\% dentro del régimen válido ($h \geq h_{\min}$) para las configuraciones óptimas por topología.

    \item \textbf{La aceleración respecto a VQE convencional es de 29--500$\times$} (Tabla~\ref{tab:cross_topo}, Figura~\ref{fig:gap_widens}). La MPNN ---siguiendo el paradigma de predicción de parámetros propuesto por \citet{miao2024} y \citet{zhang2025}--- predice parámetros en una única pasada forward clásica, eliminando el bucle de optimización iterativa. Esta eficacia se atribuye a la calidad de los datos de entrenamiento (warm-start \citep{puig2025}) y la expresividad de GINConv \citep{xu2019}.

    \item \textbf{El framework es extensible a variantes del Hamiltoniano de Ising} cuando los términos adicionales se mapean a compuertas de un qubit (TFIM + longitudinal: tasa de aprobación 90\%+, 5 topologías, sin compuertas 2Q adicionales, Tabla~\ref{tab:tfim_long}). El límite se alcanza con modelos que requieren interacciones de dos cuerpos adicionales (Heisenberg: fidelidad $\approx 0\%$ hasta $p = 4$, verificado en 46 runs independientes, Sección~\ref{sec:extensibilidad}), consistente con \citet{wiersema2020}.

    \item \textbf{La predicción funciona con alta eficiencia de datos y escala con $N$:} 16--39 puntos VQE bastan para 100\% PassRate en interpolación a $N = 30$--$60$. La variante UnifiedMPNN con parametrización bond-resolved permite entrenar un único modelo sobre múltiples $N$ simultáneamente, alcanzando 98--100\% PassRate en 4 de 5 topologías (Tabla~\ref{tab:cross_n_transfer}). El error per-site $|\Delta E|/N$ se mantiene constante al crecer $N$, confirmando escalado extensivo.

    \item \textbf{Los límites del pipeline son físicos, no algorítmicos:} la frontera de expresividad del HVA en $h \approx 1.3$ ($N = 10$) es un resultado confirmado independientemente por \citet{tripathi2026} y explicable por la cota de profundidad de \citet{sumeet2025}. Profundidad suficiente ($p \approx N-1$) elimina esta frontera completamente.
\end{enumerate}

\section{Cumplimiento de Objetivos}

\begin{itemize}
    \item \textbf{OE1} (Análisis estado del arte): Completado --- Capítulo 2, 40+ referencias 2022--2026.
    \item \textbf{OE2} (Ground truth): Completado --- Diagonalización exacta ($N \leq 14$) y DMRG ($N > 14$, $\chi = 64$--$200$), 5 topologías.
    \item \textbf{OE3} (VQE warm-start): Completado --- Fidelidad media $> 0.95$, 93\%+ de puntos $\geq 0.93$ en régimen válido.
    \item \textbf{OE4} (Predictor MPNN): Completado --- $\Delta E/\text{gap} < 5\%$ en 5 topologías dentro del régimen válido; cross-N resuelto mediante UnifiedMPNN bond-resolved (Tabla~\ref{tab:cross_n_transfer}), con escalado extensivo confirmado hasta $N = 40$ (Tabla~\ref{tab:large_n_chain}).
    \item \textbf{OE5} (Validación sistemática): Completado --- 400+ ejecuciones, 5 topologías, $N = 4$--$20$, $p = 1$--$4$.
    \item \textbf{OE6} (Extensibilidad): Completado --- TFIM longitudinal ($g \leq 0.5$, tasa de aprobación 90\%+), Heisenberg transversal (resultado negativo exhaustivo: 0\% en 26 runs), TFIM frustrado (simulación viable).
\end{itemize}

Los objetivos OE1--OE6 se cumplen íntegramente. El despliegue en hardware cuántico real queda como línea de trabajo futuro inmediata, para la cual el pipeline está preparado arquitectónicamente.

\section{Líneas de Trabajo Futuro}

El trabajo desarrollado abre las siguientes perspectivas de futuro para el campo de la simulación cuántica variacional:

\begin{enumerate}
    \item \textbf{Despliegue en hardware cuántico con mitigación de errores:} El pipeline está arquitectónicamente preparado para ejecución en hardware IBM (módulo \texttt{HardwareBackend} implementado). La combinación con técnicas de mitigación PEA-ZNE \citep{kim2023} y Dynamic Decoupling \citep{pokharel2025} permitiría validar la robustez de las predicciones MPNN frente a ruido real, extendiendo los resultados de simulación ideal presentados en este trabajo.

    \item \textbf{Extensión a parámetros bond-resolved para sistemas 2D:} Para redes bidimensionales (Kagomé, cuadrada), el HVA con parámetros por enlace (79+ dimensiones a $N = 40$) ha demostrado viabilidad intra-N en simulación (§\ref{subsec:bond_resolved}). La arquitectura MPNN con pooling global está diseñada para esta escalación, abriendo la posibilidad de abordar sistemas frustrados bidimensionales donde DMRG pierde eficiencia \citep{martin2026}.

    \item \textbf{Redes de Kagomé y frustración geométrica:} \citet{ahsan2025} demostraron viabilidad a 103 sitios en IBM Heron. La extensión del pipeline a Kagomé requiere DMRG 2D como fuente de ground truth pero no modificaciones arquitectónicas, permitiendo abordar sistemas donde el estado fundamental es genuinamente desconocido.

    \item \textbf{Técnicas generativas para cuantificación de incertidumbre:} Incorporar flujos normalizantes \citep{zou2026} para generar distribuciones de parámetros en lugar de predicciones puntuales, proporcionando estimaciones de incertidumbre intrínsecas que complementen el MC-Dropout \citep{gal2016} implementado.

    \item \textbf{Escalamiento al régimen de utilidad cuántica:} Los resultados de simulación ideal sugieren que el pipeline escala favorablemente con $N$ (precisión mejora al alejarse del punto crítico). La validación en el régimen $N > 22$ ---donde los simuladores exactos de vector de estado fallan--- mediante backends MPS y posterior ejecución en hardware constituye la extensión natural del trabajo.

    \item \textbf{Aplicación a problemas abiertos en materia condensada:} El pipeline está diseñado para escalar hacia sistemas cuyo diagrama de fases es genuinamente desconocido o computacionalmente inaccesible por métodos clásicos. Las direcciones concretas incluyen: (a) antiferromagnetos frustrados en redes de Kagomé ($N > 36$), donde el Problema del Signo impide Monte Carlo y DMRG 2D pierde fiabilidad \citep{troyer2005, ahsan2025}; (b) integración con Numerical Linked-Cluster Expansion (NLCE) \citep{sumeet2025}, donde el pipeline actuaría como solucionador de clusters, permitiendo acceso al límite termodinámico desde VQE de sistemas pequeños; y (c) detección de transiciones topológicas mediante análisis del espacio de pesos de la red entrenada \citep{hernandes2025}, proporcionando indicadores de cambio de fase sin costo cuántico adicional.
\end{enumerate}

En síntesis, la aportación de este trabajo establece un marco metodológico reproducible y validado en simulación ideal que puede ser adoptado por cualquier grupo de investigación, proporcionando una aceleración de $29$--$500\times$ en la caracterización de fases cuánticas de sistemas de espines y sentando las bases para su despliegue futuro en hardware cuántico real.


%---------------------------
% BIBLIOGRAFÍA
%---------------------------

\begin{thebibliography}{99}

\bibitem[Ahsan et al.(2025)]{ahsan2025} Ahsan, M. et al. (2025). \emph{Utility-scale quantum computation of ground-state energy in a 100+ site planar Kagome antiferromagnet via Hamiltonian engineering}. arXiv preprint arXiv:2507.06361.

\bibitem[Anderson(1972)]{anderson1972} Anderson, P. W. (1972). \emph{More is different}. Science, 177(4047), 393--396.

\bibitem[Cerezo et al.(2021)]{cerezo2021} Cerezo, M., Sone, A., Volkoff, T., Cincio, L., \& Coles, P. J. (2021). \emph{Cost function dependent barren plateaus in shallow parametrized quantum circuits}. Nature Communications, 12(1), 1791.

\bibitem[Dutta et al.(2015)]{dutta2015} Dutta, A., Aeppli, G., Chakrabarti, B. K., Divakaran, U., Rosenbaum, T. F., \& Sen, D. (2015). \emph{Quantum phase transitions in transverse field spin models: From statistical physics to quantum information}. Cambridge University Press.

\bibitem[Feynman(1982)]{feynman1982} Feynman, R. P. (1982). \emph{Simulating physics with computers}. International Journal of Theoretical Physics, 21(6/7), 467--488.

\bibitem[Fey \& Lenssen(2019)]{fey2019} Fey, M., \& Lenssen, J. E. (2019). \emph{Fast graph representation learning with PyTorch Geometric}. ICLR Workshop on Representation Learning on Graphs and Manifolds.

\bibitem[Gal \& Ghahramani(2016)]{gal2016} Gal, Y., \& Ghahramani, Z. (2016). \emph{Dropout as a Bayesian approximation: Representing model uncertainty in deep learning}. Proceedings of the 33rd ICML, 48, 1050--1059.

\bibitem[Gilmer et al.(2017)]{gilmer2017} Gilmer, J., Schoenholz, S. S., Riley, P. F., Vinyals, O., \& Dahl, G. E. (2017). \emph{Neural message passing for quantum chemistry}. Proceedings of the 34th ICML, 70, 1263--1272.

\bibitem[Hauschild \& Pollmann(2018)]{hauschild2018} Hauschild, J., \& Pollmann, F. (2018). \emph{Efficient numerical simulations with Tensor Networks: Tensor Network Python (TeNPy)}. SciPost Physics Lecture Notes, 5.

\bibitem[Hernandes et al.(2025)]{hernandes2025} Hernandes, V. et al. (2025). \emph{Adiabatic fine-tuning of neural quantum states enables detection of phase transitions in weight space}. arXiv preprint arXiv:2503.17140v2.

\bibitem[Huang et al.(2022)]{huang2022} Huang, H.-Y., Kueng, R., Torlai, G., Albert, V. V., \& Preskill, J. (2022). \emph{Provably efficient machine learning for quantum many-body problems}. Science, 377(6613), eabk3333.

\bibitem[Huang et al.(2026b)]{huang2026frustration} Huang, C. et al. (2026). \emph{Frustration-induced expressibility limitations in variational quantum algorithms}. arXiv preprint arXiv:2604.11688.

\bibitem[Kiiamov et al.(2026)]{kiiamov2026} Kiiamov, A. G. et al. (2026). \emph{Simulating Wigner localisation with the IBM Heron 2 quantum processor}. arXiv preprint arXiv:2601.01263.

\bibitem[Kim et al.(2023)]{kim2023} Kim, Y. et al. (2023). \emph{Evidence for the utility of quantum computing before fault tolerance}. Nature, 618, 500--505.

\bibitem[Kochkov et al.(2021)]{kochkov2021} Kochkov, D., Pfaff, T., Sanchez-Gonzalez, A., Battaglia, P., \& Clark, B. K. (2021). \emph{Learning ground states of quantum Hamiltonians with graph networks}. arXiv preprint arXiv:2110.06390.

\bibitem[Kumar et al.(2026)]{kumar2026} Kumar, S. et al. (2026). \emph{A study of entanglement and ansatz expressivity for the transverse-field Ising model using variational quantum eigensolver}. arXiv preprint arXiv:2602.17662.

\bibitem[Lee et al.(2026)]{lee2026} Lee, J. et al. (2026). \emph{Improving generalization and trainability of quantum eigensolvers via graph neural encoding}. arXiv preprint arXiv:2602.19752.

\bibitem[Ma et al.(2025)]{ma2025} Ma, Y. et al. (2025). \emph{Variational quantum eigensolver with zero-noise extrapolation on a superconducting processor}. arXiv preprint arXiv:2504.06554.

\bibitem[Martin et al.(2026)]{martin2026} Martin, E. et al. (2026). \emph{Tensor network methods vs.\ quantum computing: resources and boundaries for quantum advantage}. arXiv preprint arXiv:2602.04676.

\bibitem[McClean et al.(2018)]{mcclean2018} McClean, J. R., Boixo, S., Smelyanskiy, V. N., Babbush, R., \& Neven, H. (2018). \emph{Barren plateaus in quantum neural network training landscapes}. Nature Communications, 9(1), 4812.

\bibitem[Mele et al.(2022)]{mele2022} Mele, A. A., Mbeng, G. B., Santoro, G. E., Collura, M., \& Torta, P. (2022). \emph{Avoiding barren plateaus via transferability of smooth solutions in a Hamiltonian variational ansatz}. Physical Review A, 106, L060401.

\bibitem[Mele et al.(2026)]{mele2026} Mele, A. A., Angrisani, A., Ghosh, S., Khatri, S., Eisert, J., Stilck Fran\c{c}a, D., \& Quek, Y. (2026). \emph{Noise-induced shallow circuits and the absence of barren plateaus}. Nature Physics.

\bibitem[Meng et al.(2025)]{meng2025} Meng, F. et al. (2025). \emph{Output prediction of quantum circuits based on graph neural networks}. arXiv preprint arXiv:2504.00464.

\bibitem[Miao et al.(2024)]{miao2024} Miao, J., Hsieh, C.-Y., \& Zhang, S.-X. (2024). \emph{Neural-network-encoded variational quantum algorithms}. Physical Review Applied, 21, 014053.

\bibitem[Peruzzo et al.(2014)]{peruzzo2014} Peruzzo, A., McClean, J., Shadbolt, P., Yung, M.-H., Zhou, X.-Q., Love, P. J., Aspuru-Guzik, A., \& O'Brien, J. L. (2014). \emph{A variational eigenvalue solver on a photonic quantum processor}. Nature Communications, 5, 4213.

\bibitem[Pokharel et al.(2025)]{pokharel2025} Pokharel, B. et al. (2025). \emph{Empirical learning of dynamical decoupling on quantum processors}. arXiv preprint arXiv:2403.02294v2.

\bibitem[Preskill(2018)]{preskill2018} Preskill, J. (2018). \emph{Quantum computing in the NISQ era and beyond}. Quantum, 2, 79.

\bibitem[Puig et al.(2025)]{puig2025} Puig, R., Drudis, M., Thanasilp, S., \& Holmes, Z. (2025). \emph{Variational quantum simulation: A case study for understanding warm starts}. PRX Quantum, 6, 010317.

\bibitem[Sachdev(2011)]{sachdev2011} Sachdev, S. (2011). \emph{Quantum Phase Transitions}. Cambridge University Press, 2nd edition.

\bibitem[Savary \& Balents(2016)]{savary2016} Savary, L., \& Balents, L. (2016). \emph{Quantum spin liquids: a review}. Reports on Progress in Physics, 80(1), 016502.

\bibitem[Schollw\"ock(2011)]{schollwock2011} Schollw\"ock, U. (2011). \emph{The density-matrix renormalization group in the age of matrix product states}. Annals of Physics, 326(1), 96--192.

\bibitem[Schiffer et al.(2026)]{schiffer2026} Schiffer, B. et al. (2026). \emph{Scalable, self-verifying variational quantum eigensolver using adiabatic warm starts}. arXiv preprint arXiv:2602.17612.

\bibitem[Sharma(2026)]{sharma2026} Sharma, R. (2026). \emph{Quantum phase transitions in the transverse-field Ising model: A comparative study of exact, variational, and hardware-based approaches}. arXiv preprint arXiv:2601.17515.

\bibitem[Singh et al.(2025)]{singh2025} Singh, M. et al. (2025). \emph{Statistical benchmarking of optimization methods for variational quantum eigensolver under quantum noise}. arXiv preprint arXiv:2510.08727.

\bibitem[Skogh et al.(2023)]{skogh2023} Skogh, M., Leinonen, O., Lolur, P., \& Rahm, M. (2023). \emph{Accelerating variational quantum eigensolver convergence using parameter transfer}. Electronic Structure, 5, 035002.

\bibitem[Slavin(2025)]{slavin2025} Slavin, V. (2025). \emph{Graph neural network approach to predicting magnetization in quasi-one-dimensional Ising systems}. arXiv preprint arXiv:2507.17509.

\bibitem[Sumeet et al.(2025)]{sumeet2025} Sumeet, K. et al. (2025). \emph{Depth requirements for variational quantum circuits in many-body ground state preparation}. arXiv preprint arXiv:2310.07600.

\bibitem[Sun et al.(2025)]{sun2025} Sun, W. et al. (2025). \emph{Noise-mitigated variational quantum eigensolver with pre-training and zero-noise extrapolation}. arXiv preprint arXiv:2501.01646.

\bibitem[Teoh et al.(2025)]{teoh2025} Teoh, Y. H. et al. (2025). \emph{Variational quantum simulations of a two-dimensional frustrated transverse-field Ising model on a trapped-ion quantum computer}. arXiv preprint arXiv:2505.22932.

\bibitem[Tilly et al.(2022)]{tilly2022} Tilly, J. et al. (2022). \emph{The variational quantum eigensolver: A review of methods and best practices}. Physics Reports, 986, 1--128.

\bibitem[Tripathi et al.(2026)]{tripathi2026} Tripathi, A. P., Mathur, N., \& Tripathi, V. (2026). \emph{Ans\"atz expressivity and optimization in variational quantum simulations of transverse-field Ising model across system sizes}. arXiv preprint arXiv:2604.20961.

\bibitem[Troyer \& Wiese(2005)]{troyer2005} Troyer, M., \& Wiese, U. J. (2005). \emph{Computational complexity and fundamental limitations to fermionic quantum Monte Carlo simulations}. Physical Review Letters, 94(17), 170201.

\bibitem[Tsubouchi et al.(2023)]{tsubouchi2023} Tsubouchi, K., Sagawa, T., \& Takagi, K. (2023). \emph{Universal cost bound of quantum error mitigation based on quantum estimation theory}. Physical Review Letters, 131, 210602.

\bibitem[Uvarov et al.(2024)]{uvarov2024} Uvarov, A. et al. (2024). \emph{Mitigating quantum gate errors for variational eigensolvers using hardware-inspired zero-noise extrapolation}. arXiv preprint arXiv:2307.11156v3.

\bibitem[van den Berg et al.(2022)]{vandenberg2022} van den Berg, E., Minev, Z. K., Kandala, A., \& Temme, K. (2022). \emph{Probabilistic error cancellation with multi-copy Clifford circuits}. Nature Physics, 19, 752--759.

\bibitem[Wang et al.(2024)]{wang2024} Wang, S. et al. (2024). \emph{Error-mitigated variational quantum eigensolver with affine correction}. arXiv preprint arXiv:2604.16815.

\bibitem[Wiersema et al.(2020)]{wiersema2020} Wiersema, R., Zhou, C., de Sereville, Y., Carrasquilla, J., \& Kim, Y. B. (2020). \emph{Exploring entanglement and optimization within the Hamiltonian Variational Ansatz}. PRX Quantum, 1(2), 020319.

\bibitem[Xu et al.(2019)]{xu2019} Xu, K., Hu, W., Leskovec, J., \& Jegelka, S. (2019). \emph{How powerful are graph neural networks?}. Proceedings of the 7th ICLR.

\bibitem[Zhang et al.(2025)]{zhang2025} Zhang, C., Jiang, L., \& Chen, F. (2025). \emph{Qracle: A graph-neural-network-based parameter initializer for variational quantum eigensolvers}. arXiv preprint arXiv:2505.01236.

\bibitem[Zhang et al.(2025b)]{zhang2025vqezy} Zhang, C. et al. (2025). \emph{VQEzy: An open-source dataset for parameter initialization in variational quantum eigensolvers}. arXiv preprint arXiv:2509.17322.

\bibitem[Zou et al.(2026)]{zou2026} Zou, H., Rahm, M., Kockum, A. F., \& Olsson, S. (2026). \emph{Generative flow-based warm start of the variational quantum eigensolver}. npj Quantum Information, 12, 5.

\bibitem[Huang et al.(2025)]{huang2025palqo} Huang, Y., Hao, Y., Zhou, J., Yuan, X., Wang, X., \& Du, Y. (2025). \emph{PALQO: Physics-informed model for accelerating large-scale quantum optimization}. Advances in Neural Information Processing Systems (NeurIPS 2025). arXiv preprint arXiv:2509.20733.

\bibitem[Karim et al.(2025)]{karim2025} Karim, A., Zhang, S., \& Usman, M. (2025). \emph{Fast and noise-aware machine learning variational quantum eigensolver optimiser}. arXiv preprint arXiv:2503.20210.

\bibitem[Yang et al.(2024)]{yang2024pvls} Yang, Y., Zhang, C., \& Jiang, L. (2024). \emph{A learning-based parameter prediction technique for variational quantum linear solvers}. arXiv preprint arXiv:2512.04909.

\bibitem[Yin et al.(2025)]{yin2025} Yin, Z.-Q. et al. (2025). \emph{Learning variational quantum circuit parameters with classical artificial intelligence for quantum phase transition detection}. arXiv preprint arXiv:2506.06678.

\end{thebibliography}


%---------------------------
% ANEXOS
%---------------------------

\appendix

\chapter{Resultados y Mediciones Detalladas}
\label{app:resultados}

Este apéndice incluye los resultados detallados de todas las pruebas sistemáticas realizadas, tablas extendidas de mediciones, y representaciones visuales de los datos que complementan los resultados presentados en el Capítulo~5.

\section{Tablas Completas Cross-Topology ($N = 10$, $p = 2$)}

\begin{table}[h]
\centering
\caption{Resultados detallados cadena 1D, $N = 10$, $p = 2$ (top 10 variantes). Fuente: \texttt{results/thesis/}.}
\label{tab:app_chain}
\begin{tabular}{lccccc}
\toprule
\textbf{Variante} & \textbf{$\Delta E/\text{gap}$} & \textbf{Conv\%} & \textbf{$\theta$-smooth} & \textbf{Gen.gap} & \textbf{Tiempo} \\
\midrule
nl\_chain\_baseline & 0.0010 & 100\% & 0.039 & 1.7e-05 & 22s \\
nl\_chain\_1d\_baseline & 0.0022 & 100\% & 0.866 & 5.7e-05 & 22s \\
nl\_htest\_multi\_spread & 0.0148 & 100\% & 0.015 & 6.6e-04 & 39s \\
nl\_grid\_dense16 & 0.0273 & 100\% & 0.009 & 1.5e-05 & 52s \\
nl\_htest\_single\_safe & 0.0274 & 100\% & 0.015 & 2.5e-05 & 40s \\
nl\_patience\_500 & 0.0276 & 100\% & 0.033 & 4.4e-05 & 24s \\
nl\_hidden\_128 & 0.0278 & 100\% & 0.033 & 3.6e-05 & 33s \\
nl\_restarts\_5 & 0.0280 & 100\% & 0.033 & 9.7e-05 & 23s \\
nl\_seed\_42 & 0.0286 & 100\% & 0.033 & 1.7e-04 & 23s \\
nl\_seed\_44 & 0.0282 & 100\% & 0.033 & 4.8e-05 & 24s \\
\bottomrule
\end{tabular}
\end{table}

\begin{table}[h]
\centering
\caption{Resultados detallados escalera, $N = 10$, $p = 2$ (top 10 variantes).}
\label{tab:app_ladder}
\begin{tabular}{lccccc}
\toprule
\textbf{Variante} & \textbf{$\Delta E/\text{gap}$} & \textbf{Conv\%} & \textbf{$\theta$-smooth} & \textbf{Gen.gap} & \textbf{Tiempo} \\
\midrule
nl\_htest\_multi & 0.0017 & 100\% & 0.020 & 4.0e-05 & 28s \\
nl\_grid\_dense11 & 0.0164 & 100\% & 0.019 & 6.6e-06 & 72s \\
nl\_grid\_dense9 & 0.0164 & 100\% & 0.019 & 8.1e-06 & 34s \\
nl\_restarts\_1 & 0.0165 & 100\% & 0.036 & 3.1e-05 & 37s \\
nl\_seed\_43 & 0.0166 & 100\% & 0.036 & 3.9e-05 & 21s \\
nl\_hidden\_64 & 0.0168 & 100\% & 0.036 & 2.2e-05 & 20s \\
nl\_restarts\_5 & 0.0170 & 100\% & 0.036 & 2.0e-05 & 36s \\
nl\_hidden\_128 & 0.0170 & 100\% & 0.036 & 4.3e-05 & 22s \\
nl\_periodic & 0.0198 & 100\% & 0.024 & 6.7e-05 & 21s \\
nl\_seed\_44 & 0.0297 & 100\% & 0.033 & 5.7e-05 & 21s \\
\bottomrule
\end{tabular}
\end{table}

\begin{table}[h]
\centering
\caption{Resultados detallados triangular, $N = 10$, $p = 2$ (top 10 variantes).}
\label{tab:app_tri}
\begin{tabular}{lccccc}
\toprule
\textbf{Variante} & \textbf{$\Delta E/\text{gap}$} & \textbf{Conv\%} & \textbf{$\theta$-smooth} & \textbf{Gen.gap} & \textbf{Tiempo} \\
\midrule
nl\_htest\_multi & 0.0044 & 86\% & 0.014 & 1.2e-05 & 31s \\
ext\_high\_h & 0.0081 & 100\% & 0.021 & 2.6e-04 & 40s \\
nl\_restarts\_5 & 0.0371 & 100\% & 0.021 & 4.2e-05 & 25s \\
nl\_grid\_standard7 & 0.0371 & 100\% & 0.014 & 2.1e-05 & 33s \\
nl\_seed\_44 & 0.0372 & 100\% & 0.021 & 6.1e-05 & 27s \\
nl\_grid\_sparse5 & 0.0372 & 100\% & 0.021 & 5.1e-05 & 28s \\
nl\_grid\_dense9 & 0.0373 & 89\% & 0.011 & 1.0e-05 & 42s \\
nl\_hidden\_128 & 0.0376 & 100\% & 0.021 & 7.6e-05 & 32s \\
nl\_hidden\_64 & 0.0384 & 100\% & 0.021 & 5.1e-05 & 38s \\
nl\_hidden\_256 & 0.0389 & 100\% & 0.021 & 8.8e-04 & 36s \\
\bottomrule
\end{tabular}
\end{table}


\section{Escalamiento MPS Detallado ($N = 30$--$60$, Phase 3 MPNN)}

\begin{table}[h]
\centering
\caption{Phase 3 MPNN: deploy intra-N ($p = 1$, cadena 1D, backend MPS determinista, $\chi = 64$, L-BFGS-B con restarts adaptativos).}
\label{tab:app_phase3}
\begin{tabular}{rcccccc}
\toprule
\textbf{$N$} & \textbf{Pts.\ entren.} & \textbf{MSE} & \textbf{PassRate} & \textbf{$\Delta E/\text{gap}$ (media)} & \textbf{$|\Delta E|$ (media)} & \textbf{$|\Delta E|/N$} \\
\midrule
30 & 25 & $4.0 \times 10^{-6}$ & 10/10 & 1.64\% & 0.068 & 0.0023 \\
40 & 16 & $4.8 \times 10^{-6}$ & 15/15 & 1.43\% & 0.066 & 0.0016 \\
50 & 28 & $1.8 \times 10^{-5}$ & 11/11 & 0.84\% & 0.064 & 0.0013 \\
60 & 39 & $4.8 \times 10^{-6}$ & 16/16 & 0.81\% & 0.051 & 0.0009 \\
\bottomrule
\end{tabular}
\end{table}

\begin{table}[h]
\centering
\caption{Frontera de expresividad empírica por configuración ($\Delta E/\text{gap} = 5\%$, backend MPS determinista, $\chi = 64$, mediana de 3 semillas). Los ajustes lineales (Ec.~\ref{eq:frontier_p1}--\ref{eq:frontier_p2}) describen correctamente el comportamiento observado.}
\label{tab:app_frontier}
\begin{tabular}{rcccc}
\toprule
\textbf{$N$} & \textbf{$p$} & \textbf{$h_{\min}$ empírico} & \textbf{$h_{\min}$ (fit lineal)} & \textbf{Residuo} \\
\midrule
20 & 1 & 2.22 & 2.50 & $+$0.28 \\
40 & 1 & 2.61 & 2.65 & $+$0.04 \\
60 & 1 & 2.87 & 2.80 & $-$0.07 \\
80 & 1 & 3.07 & 2.94 & $-$0.13 \\
20 & 2 & 1.62 & 1.67 & $+$0.05 \\
40 & 2 & 1.80 & 1.77 & $-$0.03 \\
80 & 2 & 2.00 & 1.97 & $-$0.03 \\
20 & 3 & 1.42 & \multicolumn{2}{c}{$\approx 1.6 \pm 0.1$ (constante)} \\
80 & 3 & 1.67 & \multicolumn{2}{c}{$\approx 1.6 \pm 0.1$ (constante)} \\
\bottomrule
\end{tabular}
\end{table}

\section{Heisenberg XXZ: Escalamiento con Profundidad}

\begin{table}[h]
\centering
\caption{Escalamiento de fidelidad del HVA con profundidad para Heisenberg ($\Delta = 1.0$, $N = 6$, $h = 3.0$). TFIM como control.}
\label{tab:app_heisenberg_depth}
\begin{tabular}{lcccc}
\toprule
\textbf{$p$} & \textbf{Modelo} & \textbf{Fidelidad} & \textbf{$E_{\text{VQE}}$} & \textbf{Gap a GS} \\
\midrule
1 & Heisenberg & 0.0000 & $-$5.60 & 8.87 \\
2 & Heisenberg & 0.0020 & $-$8.60 & 5.86 \\
3 & Heisenberg & 0.3708 & $-$10.78 & 3.68 \\
5 & Heisenberg & 0.4772 & $-$13.06 & 1.40 \\
6 & Heisenberg & 0.4291 & $-$12.50 & 1.97 \\
\midrule
2 & TFIM (control) & 0.9957 & $-$9.83 & 0.02 \\
\bottomrule
\end{tabular}
\end{table}

\section{Resultados $p = 1$ Multi-Topología ($N = 10$)}

\begin{table}[h]
\centering
\caption{Pipeline $p = 1$, $N = 10$, resultados por semilla y topología.}
\label{tab:app_p1_detail}
\begin{tabular}{llcccc}
\toprule
\textbf{Topología} & \textbf{Seed} & \textbf{$h_{\text{test}}$} & \textbf{$\Delta E/\text{gap}$} & \textbf{Veredicto} & \textbf{$\theta$-smooth} \\
\midrule
chain\_1d & 42 & 2.75 & 0.042 & PASS & 0.021 \\
chain\_1d & 43 & 2.75 & 0.041 & PASS & 0.021 \\
chain\_1d & 44 & 2.75 & 0.008 & PASS & 0.021 \\
ladder & 42 & 3.25 & 0.033 & PASS & 0.014 \\
ladder & 43 & 3.25 & 0.025 & PASS & 0.014 \\
ladder & 44 & 3.25 & 0.029 & PASS & 0.014 \\
triangular & 42 & 4.25 & 0.032 & PASS & 0.011 \\
triangular & 43 & 4.25 & 0.035 & PASS & 0.011 \\
triangular & 44 & 4.25 & 0.033 & PASS & 0.011 \\
heavy-hex & 42 & 3.25 & 0.006 & PASS & 0.014 \\
heavy-hex & 43 & 3.25 & 0.006 & PASS & 0.014 \\
heavy-hex & 44 & 3.25 & 0.006 & PASS & 0.014 \\
\bottomrule
\end{tabular}
\end{table}

%---------------------------
% ANEXO B: ARTÍCULO DE INVESTIGACIÓN
%---------------------------

\chapter{Artículo de Investigación}
\label{app:articulo}

% ============================================================
% RESEARCH ARTICLE — 6-8 pages
% ============================================================

\begin{center}
{\LARGE \textbf{GNN-HVA: A Hybrid Framework for Quantum Phase Classification via Graph Neural Network Parameter Prediction on Shallow Variational Circuits}}

\vspace{0.5cm}

{\large Franco Raineri}

\vspace{0.2cm}

{\normalsize Universidad Internacional de La Rioja (UNIR), Logroño, Spain}

{\normalsize \texttt{franco.raineri@unir.net}}

\vspace{0.5cm}

{\textbf{Abstract}}
\end{center}

\noindent We present GNN-HVA, a hybrid classical-quantum pipeline for quantum phase classification. The framework combines a Message-Passing Neural Network (MPNN) based on GINConv with a Hamiltonian Variational Ansatz (HVA) to classify quantum phases of the Transverse Field Ising Model, eliminating the iterative quantum optimization loop. The MPNN predicts variational parameters directly from the Hamiltonian's graph representation, achieving a $29$--$500\times$ speedup compared to randomly-initialized VQE. Systematic validation over 400+ ideal-simulation pipeline executions across 5 lattice topologies (chain, ladder, triangular, square, heavy-hex), circuit depths $p = 1$--$4$, and system sizes $N = 4$--$20$ demonstrates $\Delta E/\text{gap} < 5\%$ within the valid operating regime. The UnifiedMPNN bond-resolved variant enables cross-$N$ generalization, maintaining extensive scaling ($|\Delta E|/N \approx$ constant) up to $N = 40$ with data trained on $N \leq 20$. We identify the HVA expressibility boundary as a physics limit quasi-independent of $N$, independently confirmed by Tripathi et al.\ (2026).

\vspace{0.3cm}
\noindent\textbf{Keywords:} VQE, Graph Neural Networks, quantum phase transitions, Ising model, variational circuits.

\section*{1. Introduction}

Variational Quantum Eigensolvers (VQE) face three fundamental barriers: barren plateaus in deep circuits \citep{mcclean2018}, noise-induced depth truncation to $O(\log n)$ layers on hardware \citep{mele2026}, and prohibitive evaluation cost ($\sim$500--1000 circuit evaluations per parameter point \citep{tilly2022}). Recent theoretical advances show that the Hamiltonian Variational Ansatz (HVA) \citep{wiersema2020} combined with warm-start optimization \citep{puig2025} produces smooth, learnable parameter landscapes.

We exploit this by training a GINConv-based MPNN \citep{xu2019} to predict optimal HVA parameters directly from the Hamiltonian's graph structure, following the parameter prediction paradigm proposed by \citet{miao2024} and \citet{zhang2025}. Unlike prior work that reduces optimization iterations (Qracle: $-$64\% \citep{zhang2025}; Flow-VQE: $50\times$ \citep{zou2026}), our pipeline achieves sub-5\% accuracy without iterative refinement, attributed to the high-quality training data from warm-start VQE \citep{puig2025, schiffer2026}. All results are validated via ideal simulation (StatevectorEstimator), establishing a theoretical performance ceiling for future hardware deployment.

\section*{2. Method}

The GNN-HVA pipeline operates in three phases:
\begin{enumerate}
    \item \textbf{Phase 1:} Ground truth generation via exact diagonalization ($N \leq 14$) or DMRG ($N > 14$, $\chi = 64$) \citep{hauschild2018}.
    \item \textbf{Phase 2:} Warm-start VQE on HVA ($p = 1$--$4$) with descending sweep $h_{\max} \to h_{\min}$, producing high-quality $\theta^*(h)$ training data \citep{puig2025}.
    \item \textbf{Phase 3:} MPNN training (3-layer GINConv + global\_mean\_pool + MLP) to learn $h \mapsto \theta^*$ from graph representation \citep{fey2019}, followed by deployment evaluation.
\end{enumerate}

The TFIM Hamiltonian $H = -J\sum Z_iZ_{i+1} - h\sum X_i$ is represented as a graph where nodes carry the field attribute $h_i$ and edges carry the coupling $J_{ij}$. The success criterion is $\Delta E/\text{gap} < 5\%$ with correct phase classification via local observables.

\section*{3. Results}

\textbf{Cross-topology validation ($N = 10$):} The pipeline achieves 95\% pass rate on chain\_1d ($p = 3$), 90\% on heavy-hex ($p = 3$), and 56--79\% on 2D topologies ($p = 4$). Topology dominates over circuit depth: 2D topologies gain only +4--10\% from $p = 2 \to 4$, while chain/heavy-hex gain +19--20\%.

\textbf{Scaling ($N = 4$--$20$):} Accuracy improves with system size within the valid regime. With restricted $h$-range ($h \geq 1.3$), chain\_1d achieves 100\% pass rate at $N = 10$ ($p = 3$), $N = 16$ ($p = 4$), and $N = 20$ ($p = 4$). Speedup: $29$--$500\times$ over random initialization. UnifiedMPNN bond-resolved enables cross-$N$ generalization with extensive scaling up to $N = 40$.

\textbf{Ising variants:} TFIM + longitudinal field achieves pass rate $90\%+$ on heavy-hex ($p = 3$) with zero additional 2-qubit gates. Heisenberg transversal shows 0\% pass rate across all topologies and depths ($p = 1$--$4$, 26 runs), confirming an expressibility limit confirmed independently by \citet{wiersema2020}.

\section*{4. Discussion}

The pipeline's effectiveness rests on four converging properties: (1) smooth HVA landscape within a phase \citep{mele2022}, (2) graph structure informing predictions \citep{slavin2025}, (3) moderate circuit depth compatible with hardware constraints \citep{mele2026}, and (4) high-quality training data from warm-start \citep{puig2025}. The expressibility boundary at $h \approx 1.3$ ($N = 10$) is a physics limit independently confirmed by \citet{tripathi2026}, not a pipeline failure. Sufficient depth ($p \approx N-1$) eliminates this boundary entirely.

Compared to literature: NN-VQE \citep{miao2024} validates MLP prediction at $N \leq 10$; Qracle \citep{zhang2025} validates GNN initialization reducing 64\% of iterations. Our contribution is the integration and systematic validation across 5 topologies, 4 depth levels, and multiple Hamiltonian variants --- extending the paradigm rather than inventing it.

\section*{5. Conclusions}

GNN-HVA demonstrates that the integration of GNN-based parameter prediction \citep{miao2024, zhang2025}, HVA circuits \citep{wiersema2020}, and warm-start optimization \citep{puig2025} into a unified pipeline enables quantum phase classification with $29$--$500\times$ speedup over conventional VQE. The pipeline achieves 95--100\% pass rate on optimal configurations within the valid operating regime, extends to Ising variants, and enables cross-$N$ generalization via UnifiedMPNN bond-resolved architecture. These ideal-simulation results establish the methodology for future deployment on quantum hardware with error mitigation \citep{kim2023}. Code available at: \url{https://github.com/franraineri/qmbp_gnn_adapt-vqe}.

\vspace{0.3cm}
\noindent\textbf{Acknowledgments.} This work was developed as part of a Master's Thesis at UNIR.


\end{document}
